% !TEX program = lualatex
% ---------------------------------------------------------------------------------------------
%  Readback note 13: Glued lattices, product forms and the integral comparison forms
%
%  Build:  python3 tools/rb.py note <slug>
%
%  One unit of a Lean artifact, rewritten as ordinary mathematics. The prose is hand-written. Lean names,
%  line ranges and counts come from data.tex; quoted Lean statements are copied from the artifact by the build.
%
%  Rules of the text: standard terminology; every symbol is explained where it first appears and listed again
%  in the notation at the end; no letter has two meanings; a proof is complete or says what it takes for granted;
%  Lean names appear only in the LEAN lines and in the concordance.
%
%  The signature rule (the build enforces it): a function, relation or predicate that the artifact defines is
%  written only through its notation macro, which the build makes from the spec (`notation`) with one slot for
%  each explicit argument of the Lean definition. No argument is fixed and dropped, also not "for brevity" and
%  not for one section. Every argument and every hypothesis of a stated declaration stands inside its
%  statement, or in a Setting that the statement cites by number; only implicit and instance arguments (the
%  ambient objects) may stand in the conventions.
% ---------------------------------------------------------------------------------------------
\documentclass[11pt,a4paper]{article}
\ifdefined\pdfvariable\pdfvariable suppressoptionalinfo \numexpr 1+2+4+8\relax\fi
\usepackage[a4paper,hmargin=35mm,top=22mm,bottom=27mm,footskip=12mm,nohead]{geometry}

\usepackage{fontspec}
\directlua{
  luaotfload.add_fallback("gfallback", {"STIXTwoMath-Regular.otf:mode=node;", "DejaVuSans.ttf:mode=node;"})
  luaotfload.add_fallback("gmonofallback", {"DejaVuSansMono.ttf:mode=node;", "STIXTwoMath-Regular.otf:mode=node;"})
}
\setmainfont{NewCM10}[Extension=.otf,UprightFont=*-Book,ItalicFont=*-BookItalic,
  BoldFont=*-Bold,BoldItalicFont=*-BoldItalic,SmallCapsFont=*-Book,SmallCapsFeatures={Letters=SmallCaps},
  RawFeature={fallback=gfallback}]
\setsansfont{NewCMSans10}[Extension=.otf,UprightFont=*-Book,ItalicFont=*-BookOblique,
  BoldFont=*-Bold,BoldItalicFont=*-BoldOblique,RawFeature={fallback=gfallback}]
\setmonofont{JetBrainsMono}[Extension=.otf,UprightFont=*-Regular,ItalicFont=*-Italic,
  BoldFont=*-SemiBold,BoldItalicFont=*-SemiBoldItalic,Scale=MatchLowercase,
  RawFeature={-calt;-liga;fallback=gmonofallback}]
\usepackage{amsmath}
\usepackage{amsthm}
\usepackage[math-style=TeX]{unicode-math}
\setmathfont{NewCMMath-Book.otf}
\emergencystretch=2.5em\relax

\usepackage{xcolor}
\usepackage{tikz}
\usetikzlibrary{calc,arrows.meta,positioning,backgrounds}
\usepackage[protrusion=true,expansion=true]{microtype}
\usepackage{enumitem}
\usepackage{needspace}
\usepackage{booktabs}
\usepackage{array}
\usepackage{longtable}
\usepackage{fancyhdr}
\usepackage{lastpage}
\usepackage{multicol}
\usepackage[hidelinks,pdfusetitle]{hyperref}

\definecolor{ink}{HTML}{0F172A}
\definecolor{inksoft}{HTML}{475569}
\definecolor{inkfaint}{HTML}{94A3B8}
\definecolor{rulegrey}{HTML}{CBD5E1}
\definecolor{arith}{HTML}{C2410C}
\definecolor{est}{HTML}{7C3AED}
\definecolor{ana}{HTML}{0891B2}
\definecolor{ok}{HTML}{047857}
\definecolor{leancol}{HTML}{3730A3}

\title{Glued lattices, product forms and the integral comparison forms}
\author{grh-compact}
\date{}

% ------------------------------------------------------------------------------------- macros
\newcommand{\leanfont}{\ttfamily\fontsize{8.6}{10.5}\selectfont\color{leancol}}
\newcommand{\lean}[1]{\csname LN@#1\endcsname}
\newcommand{\fillin}[1]{{\color{arith}\itshape[#1]}}  % a place of the template that is not written yet; the build reports every one that is left
\newcommand{\leanquote}[1]{\csname LQ@#1\endcsname}
% generated by rb.py note; do not edit
\newcommand{\ConcDecls}{57}
\newcommand{\Pin}{72fead10}
\newcommand{\PinText}{that of commit {\leanfont 72fead10}}
\newcommand{\ArtifactName}{grh-compact}
\expandafter\newcommand\csname LN@sumdiagramWeight\endcsname{\href{https://github.com/lsernau/grh-lean/blob/72fead102f663d5c7a407964a13dc0cf9d977207/grh-compact/src/New/Rows/MemberKernels.lean\#L31-L35}{{\leanfont FitsRanks.\allowbreak{}sum\_diagramWeight\_eq\_boxes}}}
\expandafter\newcommand\csname LN@existsalgEquiv\endcsname{\href{https://github.com/lsernau/grh-lean/blob/72fead102f663d5c7a407964a13dc0cf9d977207/grh-compact/src/New/Rows/MemberKernels.lean\#L37-L40}{{\leanfont Parameters.\allowbreak{}exists\_algEquiv\_norm\_eq}}}
\expandafter\newcommand\csname LN@normjetle\endcsname{\href{https://github.com/lsernau/grh-lean/blob/72fead102f663d5c7a407964a13dc0cf9d977207/grh-compact/src/New/Rows/MemberKernels.lean\#L47-L88}{{\leanfont Parameters.\allowbreak{}norm\_jet\_le\_of\_isAdmissibleComb}}}
\expandafter\newcommand\csname LN@boxesle\endcsname{\href{https://github.com/lsernau/grh-lean/blob/72fead102f663d5c7a407964a13dc0cf9d977207/grh-compact/src/New/Rows/MemberKernels.lean\#L92-L96}{{\leanfont boxes\_le\_of\_isHCParam}}}
\expandafter\newcommand\csname LN@existsisHCParamofle\endcsname{\href{https://github.com/lsernau/grh-lean/blob/72fead102f663d5c7a407964a13dc0cf9d977207/grh-compact/src/New/Rows/MemberKernels.lean\#L98-L107}{{\leanfont exists\_isHCParam\_of\_le}}}
\expandafter\newcommand\csname LN@existsisHCParamofhasGap\endcsname{\href{https://github.com/lsernau/grh-lean/blob/72fead102f663d5c7a407964a13dc0cf9d977207/grh-compact/src/New/Rows/MemberKernels.lean\#L111-L113}{{\leanfont exists\_isHCParam\_of\_hasGap}}}
\expandafter\newcommand\csname LN@modKernelembFin\endcsname{\href{https://github.com/lsernau/grh-lean/blob/72fead102f663d5c7a407964a13dc0cf9d977207/grh-compact/src/New/Rows/MemberKernels.lean\#L118-L125}{{\leanfont modKernel\_embFin\_mul}}}
\expandafter\newcommand\csname LN@summemintegralAbove\endcsname{\href{https://github.com/lsernau/grh-lean/blob/72fead102f663d5c7a407964a13dc0cf9d977207/grh-compact/src/New/Rows/MemberKernels.lean\#L127-L130}{{\leanfont sum\_mem\_integralAbove}}}
\expandafter\newcommand\csname LN@existsmodKernelbound\endcsname{\href{https://github.com/lsernau/grh-lean/blob/72fead102f663d5c7a407964a13dc0cf9d977207/grh-compact/src/New/Rows/MemberKernels.lean\#L143-L167}{{\leanfont exists\_modKernel\_integral\_bound}}}
\expandafter\newcommand\csname LN@existsreduce\endcsname{\href{https://github.com/lsernau/grh-lean/blob/72fead102f663d5c7a407964a13dc0cf9d977207/grh-compact/src/New/Rows/MemberKernels.lean\#L172-L179}{{\leanfont exists\_reduceAt\_mul\_incl\_eq}}}
\expandafter\newcommand\csname LN@pointForm\endcsname{\href{https://github.com/lsernau/grh-lean/blob/72fead102f663d5c7a407964a13dc0cf9d977207/grh-compact/src/New/Rows/MemberKernels.lean\#L183-L184}{{\leanfont pointForm}}}
\expandafter\newcommand\csname LN@innerpointForm\endcsname{\href{https://github.com/lsernau/grh-lean/blob/72fead102f663d5c7a407964a13dc0cf9d977207/grh-compact/src/New/Rows/MemberKernels.lean\#L186-L191}{{\leanfont inner\_pointForm}}}
\expandafter\newcommand\csname LN@thetapointForm\endcsname{\href{https://github.com/lsernau/grh-lean/blob/72fead102f663d5c7a407964a13dc0cf9d977207/grh-compact/src/New/Rows/MemberKernels.lean\#L194-L203}{{\leanfont theta\_pointForm}}}
\expandafter\newcommand\csname LN@kernelVector\endcsname{\href{https://github.com/lsernau/grh-lean/blob/72fead102f663d5c7a407964a13dc0cf9d977207/grh-compact/src/New/Rows/MemberKernels.lean\#L207-L208}{{\leanfont kernelVector}}}
\expandafter\newcommand\csname LN@rowFixingLevel\endcsname{\href{https://github.com/lsernau/grh-lean/blob/72fead102f663d5c7a407964a13dc0cf9d977207/grh-compact/src/New/Rows/MemberKernels.lean\#L212-L214}{{\leanfont rowFixingLevel}}}
\expandafter\newcommand\csname LN@cosetKernelSum\endcsname{\href{https://github.com/lsernau/grh-lean/blob/72fead102f663d5c7a407964a13dc0cf9d977207/grh-compact/src/New/Rows/MemberKernels.lean\#L218-L221}{{\leanfont cosetKernelSum}}}
\expandafter\newcommand\csname LN@depthindicatorjet\endcsname{\href{https://github.com/lsernau/grh-lean/blob/72fead102f663d5c7a407964a13dc0cf9d977207/grh-compact/src/New/Rows/GluedFormTheta.lean\#L33-L143}{{\leanfont Data.\allowbreak{}depth\_indicator\_jet}}}
\expandafter\newcommand\csname LN@DividedLattice\endcsname{\href{https://github.com/lsernau/grh-lean/blob/72fead102f663d5c7a407964a13dc0cf9d977207/grh-compact/src/New/Rows/GluedFormTheta.lean\#L147-L151}{{\leanfont DividedLattice}}}
\expandafter\newcommand\csname LN@reducedInverse\endcsname{\href{https://github.com/lsernau/grh-lean/blob/72fead102f663d5c7a407964a13dc0cf9d977207/grh-compact/src/New/Rows/GluedFormTheta.lean\#L159-L160}{{\leanfont DividedLattice.\allowbreak{}reducedInverse}}}
\expandafter\newcommand\csname LN@weilScalar\endcsname{\href{https://github.com/lsernau/grh-lean/blob/72fead102f663d5c7a407964a13dc0cf9d977207/grh-compact/src/New/Rows/GluedFormTheta.lean\#L164-L165}{{\leanfont DividedLattice.\allowbreak{}weilScalar}}}
\expandafter\newcommand\csname LN@reduceAtconj\endcsname{\href{https://github.com/lsernau/grh-lean/blob/72fead102f663d5c7a407964a13dc0cf9d977207/grh-compact/src/New/Rows/GluedFormTheta.lean\#L172-L191}{{\leanfont DividedLattice.\allowbreak{}reduceAt\_conj\_eq\_one}}}
\expandafter\newcommand\csname LN@existsreducedInverse\endcsname{\href{https://github.com/lsernau/grh-lean/blob/72fead102f663d5c7a407964a13dc0cf9d977207/grh-compact/src/New/Rows/GluedFormTheta.lean\#L194-L208}{{\leanfont DividedLattice.\allowbreak{}exists\_reducedInverse\_eq}}}
\expandafter\newcommand\csname LN@stackOverGenScalar\endcsname{\href{https://github.com/lsernau/grh-lean/blob/72fead102f663d5c7a407964a13dc0cf9d977207/grh-compact/src/New/Rows/GluedFormTheta.lean\#L212-L213}{{\leanfont DividedSplitting.\allowbreak{}stackOverGenScalar}}}
\expandafter\newcommand\csname LN@rowFixingLevelle\endcsname{\href{https://github.com/lsernau/grh-lean/blob/72fead102f663d5c7a407964a13dc0cf9d977207/grh-compact/src/New/Rows/GluedFormTheta.lean\#L219-L222}{{\leanfont rowFixingLevel\_le\_levelOf}}}
\expandafter\newcommand\csname LN@congruenceGroup\endcsname{\href{https://github.com/lsernau/grh-lean/blob/72fead102f663d5c7a407964a13dc0cf9d977207/grh-compact/src/New/Rows/GluedFormTheta.lean\#L226-L227}{{\leanfont DividedSplitting.\allowbreak{}congruenceGroup}}}
\expandafter\newcommand\csname LN@reduceAtliftMod\endcsname{\href{https://github.com/lsernau/grh-lean/blob/72fead102f663d5c7a407964a13dc0cf9d977207/grh-compact/src/New/Rows/GluedFormTheta.lean\#L230-L232}{{\leanfont DividedSplitting.\allowbreak{}reduceAt\_liftMod\_congruenceGroup}}}
\expandafter\newcommand\csname LN@IsGlued\endcsname{\href{https://github.com/lsernau/grh-lean/blob/72fead102f663d5c7a407964a13dc0cf9d977207/grh-compact/src/New/Rows/GluedFormTheta.lean\#L236-L242}{{\leanfont DividedSplitting.\allowbreak{}IsGlued}}}
\expandafter\newcommand\csname LQ@IsGlued\endcsname{\leanquoteline{0}{structure DividedSplitting.IsGlued where}\leanquoteline{2}{indicator\_eq\_sum (x : Matrix n r O.chosenModulusSq.ResRing) (c : ℂ) : ite (L.reducedInverse * x =}\leanquoteline{4}{V.genResidue • inclMatrix \_ O.incl) c 0 = ∑ Y : V.congruenceGroup, ite}\leanquoteline{4}{((x * Y.1.1).submatrix (V.ε ∘ inr) id = inclMatrix \_ z ∧ V.stackOverGenScalar s *}\leanquoteline{4}{(x * Y.1.1).submatrix (V.ε ∘ inl) id = V.stackOverGenScalar s * inclMatrix \_ V.fixedEmb) c 0}\leanquoteline{2}{exists\_cols\_eq\_zero : ∃ K : GL N O.chosenModulusSq.ResRing, ∀ x, ∀ y ∉ univ.map V.dividedEmb,}\leanquoteline{4}{V.genResidue * (L.reducedInverse.submatrix id (V.ε ∘ inl) * K.1) x y = 0}}
\expandafter\newcommand\csname LN@productForm\endcsname{\href{https://github.com/lsernau/grh-lean/blob/72fead102f663d5c7a407964a13dc0cf9d977207/grh-compact/src/New/Rows/GluedFormTheta.lean\#L246-L248}{{\leanfont DividedSplitting.\allowbreak{}productForm}}}
\expandafter\newcommand\csname LN@existscongruenceGroupinv\endcsname{\href{https://github.com/lsernau/grh-lean/blob/72fead102f663d5c7a407964a13dc0cf9d977207/grh-compact/src/New/Rows/GluedFormTheta.lean\#L251-L256}{{\leanfont DividedSplitting.\allowbreak{}exists\_congruenceGroup\_inv\_eq\_one\_sub}}}
\expandafter\newcommand\csname LN@existsisGlued\endcsname{\href{https://github.com/lsernau/grh-lean/blob/72fead102f663d5c7a407964a13dc0cf9d977207/grh-compact/src/New/Rows/GluedFormTheta.lean\#L259-L294}{{\leanfont DividedSplitting.\allowbreak{}exists\_isGlued}}}
\expandafter\newcommand\csname LN@thetapointFormL\endcsname{\href{https://github.com/lsernau/grh-lean/blob/72fead102f663d5c7a407964a13dc0cf9d977207/grh-compact/src/New/Rows/GluedFormTheta.lean\#L308-L336}{{\leanfont DividedSplitting.\allowbreak{}theta\_pointForm}}}
\expandafter\newcommand\csname LN@thetaproductForm\endcsname{\href{https://github.com/lsernau/grh-lean/blob/72fead102f663d5c7a407964a13dc0cf9d977207/grh-compact/src/New/Rows/GluedFormTheta.lean\#L343-L366}{{\leanfont DividedSplitting.\allowbreak{}theta\_productForm}}}
\expandafter\newcommand\csname LN@thetastdClassSumproductForm\endcsname{\href{https://github.com/lsernau/grh-lean/blob/72fead102f663d5c7a407964a13dc0cf9d977207/grh-compact/src/New/Rows/GluedFormTheta.lean\#L372-L378}{{\leanfont DividedSplitting.\allowbreak{}theta\_stdClassSum\_productForm}}}
\expandafter\newcommand\csname LN@normconjintAlgPart\endcsname{\href{https://github.com/lsernau/grh-lean/blob/72fead102f663d5c7a407964a13dc0cf9d977207/grh-compact/src/New/Rows/MonomialCoefficientVector.lean\#L20-L36}{{\leanfont Data.\allowbreak{}Pair.\allowbreak{}norm\_conj\_intAlgPart\_le}}}
\expandafter\newcommand\csname LN@iotagoodProjTranslate\endcsname{\href{https://github.com/lsernau/grh-lean/blob/72fead102f663d5c7a407964a13dc0cf9d977207/grh-compact/src/New/Rows/MonomialCoefficientVector.lean\#L38-L41}{{\leanfont Data.\allowbreak{}iota\_goodProjTranslate\_apply}}}
\expandafter\newcommand\csname LN@goodProjTranslateintegral\endcsname{\href{https://github.com/lsernau/grh-lean/blob/72fead102f663d5c7a407964a13dc0cf9d977207/grh-compact/src/New/Rows/MonomialCoefficientVector.lean\#L47-L54}{{\leanfont Data.\allowbreak{}goodProjTranslate\_mem\_integralForms}}}
\expandafter\newcommand\csname LN@normSqAtgoodProjTranslate\endcsname{\href{https://github.com/lsernau/grh-lean/blob/72fead102f663d5c7a407964a13dc0cf9d977207/grh-compact/src/New/Rows/MonomialCoefficientVector.lean\#L58-L70}{{\leanfont Data.\allowbreak{}normSqAt\_goodProjTranslate\_le}}}
\expandafter\newcommand\csname LN@existsintvector\endcsname{\href{https://github.com/lsernau/grh-lean/blob/72fead102f663d5c7a407964a13dc0cf9d977207/grh-compact/src/New/Rows/MonomialCoefficientVector.lean\#L80-L94}{{\leanfont exists\_int\_vector\_harmonicPoly\_eq}}}
\expandafter\newcommand\csname LN@descentsVanishhypSum\endcsname{\href{https://github.com/lsernau/grh-lean/blob/72fead102f663d5c7a407964a13dc0cf9d977207/grh-compact/src/New/Rows/MemberOrthogonality.lean\#L35-L53}{{\leanfont descentsVanish\_hypSum}}}
\expandafter\newcommand\csname LN@stdClassSumperiod\endcsname{\href{https://github.com/lsernau/grh-lean/blob/72fead102f663d5c7a407964a13dc0cf9d977207/grh-compact/src/New/Rows/MemberOrthogonality.lean\#L62-L120}{{\leanfont stdClassSum\_period\_eq\_zero}}}
\expandafter\newcommand\csname LN@stdClassSuminnerproductForm\endcsname{\href{https://github.com/lsernau/grh-lean/blob/72fead102f663d5c7a407964a13dc0cf9d977207/grh-compact/src/New/Rows/MemberOrthogonality.lean\#L126-L138}{{\leanfont DividedSplitting.\allowbreak{}stdClassSum\_inner\_productForm\_eq\_zero}}}
\expandafter\newcommand\csname LN@loss\endcsname{\href{https://github.com/lsernau/grh-lean/blob/72fead102f663d5c7a407964a13dc0cf9d977207/grh-compact/src/New/Rows/CartanIntegrality.lean\#L19-L20}{{\leanfont DividedSplitting.\allowbreak{}loss}}}
\expandafter\newcommand\csname LN@normschurinvcols\endcsname{\href{https://github.com/lsernau/grh-lean/blob/72fead102f663d5c7a407964a13dc0cf9d977207/grh-compact/src/New/Rows/CartanIntegrality.lean\#L29-L51}{{\leanfont DividedSplitting.\allowbreak{}norm\_schur\_inv\_cols\_le\_weylGauge}}}
\expandafter\newcommand\csname LN@weylGaugecartanProduct\endcsname{\href{https://github.com/lsernau/grh-lean/blob/72fead102f663d5c7a407964a13dc0cf9d977207/grh-compact/src/New/Rows/CartanIntegrality.lean\#L56-L69}{{\leanfont DividedSplitting.\allowbreak{}weylGauge\_cartanProduct\_le}}}
\expandafter\newcommand\csname LN@smallHookNumber\endcsname{\href{https://github.com/lsernau/grh-lean/blob/72fead102f663d5c7a407964a13dc0cf9d977207/grh-compact/src/New/Rows/ComparisonForms.lean\#L17-L17}{{\leanfont smallHookNumber}}}
\expandafter\newcommand\csname LN@SmallFormValues\endcsname{\href{https://github.com/lsernau/grh-lean/blob/72fead102f663d5c7a407964a13dc0cf9d977207/grh-compact/src/New/Rows/ComparisonForms.lean\#L21-L22}{{\leanfont SmallFormValues}}}
\expandafter\newcommand\csname LN@Tables\endcsname{\href{https://github.com/lsernau/grh-lean/blob/72fead102f663d5c7a407964a13dc0cf9d977207/grh-compact/src/New/Rows/ComparisonForms.lean\#L32-L32}{{\leanfont Data.\allowbreak{}Tables}}}
\expandafter\newcommand\csname LN@splittingV\endcsname{\href{https://github.com/lsernau/grh-lean/blob/72fead102f663d5c7a407964a13dc0cf9d977207/grh-compact/src/New/Rows/ComparisonForms.lean\#L38-L42}{{\leanfont splittingV}}}
\expandafter\newcommand\csname LN@classAutrowFixingLevel\endcsname{\href{https://github.com/lsernau/grh-lean/blob/72fead102f663d5c7a407964a13dc0cf9d977207/grh-compact/src/New/Rows/ComparisonForms.lean\#L46-L49}{{\leanfont classAut\_rowFixingLevel\_eq\_bot}}}
\expandafter\newcommand\csname LN@glueLattice\endcsname{\href{https://github.com/lsernau/grh-lean/blob/72fead102f663d5c7a407964a13dc0cf9d977207/grh-compact/src/New/Rows/ComparisonForms.lean\#L53-L54}{{\leanfont glueLattice}}}
\expandafter\newcommand\csname LN@lineScale\endcsname{\href{https://github.com/lsernau/grh-lean/blob/72fead102f663d5c7a407964a13dc0cf9d977207/grh-compact/src/New/Rows/ComparisonForms.lean\#L59-L59}{{\leanfont lineScale}}}
\expandafter\newcommand\csname LN@comparisonForm\endcsname{\href{https://github.com/lsernau/grh-lean/blob/72fead102f663d5c7a407964a13dc0cf9d977207/grh-compact/src/New/Rows/ComparisonForms.lean\#L70-L75}{{\leanfont comparisonForm}}}
\expandafter\newcommand\csname LQ@comparisonForm\endcsname{\leanquoteline{0}{@[simps -isSimp] def comparisonForm : Tables j →ₗ[ℂ] (D.auxPair ϑ j).FormV where}\leanquoteline{2}{toFun ζ := (D.gProd \textasciicircum{} (splittingV j ϑ).loss : ℂ) • ∑ b, stdClassSum (rowFixingLevel \_ ((splittingV}\leanquoteline{4}{j ϑ).stackOverGenScalar b.1.1) (splittingV j ϑ).fixedEmb) fun h ↦ (splittingV j ϑ).productForm}\leanquoteline{4}{(glueLattice j ϑ b.1) (h, b.2.out) (star ((lineScale j).1 • lineMultiplier ϑ j b.1 h)) (ζ b)}\leanquoteline{2}{map\_add' := by simp [↓classAut\_rowFixingLevel\_eq\_bot]}\leanquoteline{2}{map\_smul' a := by simp [↓classAut\_rowFixingLevel\_eq\_bot, smul\_comm a, ← smul\_sum]}}
\expandafter\newcommand\csname LN@inneradjoint\endcsname{\href{https://github.com/lsernau/grh-lean/blob/72fead102f663d5c7a407964a13dc0cf9d977207/grh-compact/src/New/Rows/ComparisonForms.lean\#L81-L111}{{\leanfont inner\_adjoint\_comparisonForm}}}
\expandafter\newcommand\csname LN@comparisonFormintegral\endcsname{\href{https://github.com/lsernau/grh-lean/blob/72fead102f663d5c7a407964a13dc0cf9d977207/grh-compact/src/New/Rows/ComparisonForms.lean\#L116-L126}{{\leanfont comparisonForm\_mem\_integralForms}}}
\expandafter\newcommand\csname LN@normSqAtcomparisonForm\endcsname{\href{https://github.com/lsernau/grh-lean/blob/72fead102f663d5c7a407964a13dc0cf9d977207/grh-compact/src/New/Rows/ComparisonForms.lean\#L131-L163}{{\leanfont normSqAt\_comparisonForm\_le}}}
\expandafter\newcommand\csname LN@tableThetaadjoint\endcsname{\href{https://github.com/lsernau/grh-lean/blob/72fead102f663d5c7a407964a13dc0cf9d977207/grh-compact/src/New/Rows/ComparisonForms.lean\#L170-L192}{{\leanfont tableTheta\_adjoint\_eq\_zero}}}
\expandafter\newcommand\csname LN@existscomparisonforms\endcsname{\href{https://github.com/lsernau/grh-lean/blob/72fead102f663d5c7a407964a13dc0cf9d977207/grh-compact/src/New/Rows/ComparisonForms.lean\#L200-L227}{{\leanfont exists\_comparison\_forms}}}
\newcommand{\ConcordanceRows}{\crow{sumdiagramWeight}{\href{https://github.com/lsernau/grh-lean/blob/72fead102f663d5c7a407964a13dc0cf9d977207/grh-compact/src/New/Rows/MemberKernels.lean\#L31-L35}{{\leanfont FitsRanks.\allowbreak{}sum\_diagramWeight\_eq\_boxes}}}{MemberKernels.lean}{31--35}
\crow{existsalgEquiv}{\href{https://github.com/lsernau/grh-lean/blob/72fead102f663d5c7a407964a13dc0cf9d977207/grh-compact/src/New/Rows/MemberKernels.lean\#L37-L40}{{\leanfont Parameters.\allowbreak{}exists\_algEquiv\_norm\_eq}}}{MemberKernels.lean}{37--40}
\crow{normjetle}{\href{https://github.com/lsernau/grh-lean/blob/72fead102f663d5c7a407964a13dc0cf9d977207/grh-compact/src/New/Rows/MemberKernels.lean\#L47-L88}{{\leanfont Parameters.\allowbreak{}norm\_jet\_le\_of\_isAdmissibleComb}}}{MemberKernels.lean}{47--88}
\crow{boxesle}{\href{https://github.com/lsernau/grh-lean/blob/72fead102f663d5c7a407964a13dc0cf9d977207/grh-compact/src/New/Rows/MemberKernels.lean\#L92-L96}{{\leanfont boxes\_le\_of\_isHCParam}}}{MemberKernels.lean}{92--96}
\crow{existsisHCParamofle}{\href{https://github.com/lsernau/grh-lean/blob/72fead102f663d5c7a407964a13dc0cf9d977207/grh-compact/src/New/Rows/MemberKernels.lean\#L98-L107}{{\leanfont exists\_isHCParam\_of\_le}}}{MemberKernels.lean}{98--107}
\crow{existsisHCParamofhasGap}{\href{https://github.com/lsernau/grh-lean/blob/72fead102f663d5c7a407964a13dc0cf9d977207/grh-compact/src/New/Rows/MemberKernels.lean\#L111-L113}{{\leanfont exists\_isHCParam\_of\_hasGap}}}{MemberKernels.lean}{111--113}
\crow{modKernelembFin}{\href{https://github.com/lsernau/grh-lean/blob/72fead102f663d5c7a407964a13dc0cf9d977207/grh-compact/src/New/Rows/MemberKernels.lean\#L118-L125}{{\leanfont modKernel\_embFin\_mul}}}{MemberKernels.lean}{118--125}
\crow{summemintegralAbove}{\href{https://github.com/lsernau/grh-lean/blob/72fead102f663d5c7a407964a13dc0cf9d977207/grh-compact/src/New/Rows/MemberKernels.lean\#L127-L130}{{\leanfont sum\_mem\_integralAbove}}}{MemberKernels.lean}{127--130}
\crow{existsmodKernelbound}{\href{https://github.com/lsernau/grh-lean/blob/72fead102f663d5c7a407964a13dc0cf9d977207/grh-compact/src/New/Rows/MemberKernels.lean\#L143-L167}{{\leanfont exists\_modKernel\_integral\_bound}}}{MemberKernels.lean}{143--167}
\crow{existsreduce}{\href{https://github.com/lsernau/grh-lean/blob/72fead102f663d5c7a407964a13dc0cf9d977207/grh-compact/src/New/Rows/MemberKernels.lean\#L172-L179}{{\leanfont exists\_reduceAt\_mul\_incl\_eq}}}{MemberKernels.lean}{172--179}
\crow{pointForm}{\href{https://github.com/lsernau/grh-lean/blob/72fead102f663d5c7a407964a13dc0cf9d977207/grh-compact/src/New/Rows/MemberKernels.lean\#L183-L184}{{\leanfont pointForm}}}{MemberKernels.lean}{183--184}
\crow{innerpointForm}{\href{https://github.com/lsernau/grh-lean/blob/72fead102f663d5c7a407964a13dc0cf9d977207/grh-compact/src/New/Rows/MemberKernels.lean\#L186-L191}{{\leanfont inner\_pointForm}}}{MemberKernels.lean}{186--191}
\crow{thetapointForm}{\href{https://github.com/lsernau/grh-lean/blob/72fead102f663d5c7a407964a13dc0cf9d977207/grh-compact/src/New/Rows/MemberKernels.lean\#L194-L203}{{\leanfont theta\_pointForm}}}{MemberKernels.lean}{194--203}
\crow{kernelVector}{\href{https://github.com/lsernau/grh-lean/blob/72fead102f663d5c7a407964a13dc0cf9d977207/grh-compact/src/New/Rows/MemberKernels.lean\#L207-L208}{{\leanfont kernelVector}}}{MemberKernels.lean}{207--208}
\crow{rowFixingLevel}{\href{https://github.com/lsernau/grh-lean/blob/72fead102f663d5c7a407964a13dc0cf9d977207/grh-compact/src/New/Rows/MemberKernels.lean\#L212-L214}{{\leanfont rowFixingLevel}}}{MemberKernels.lean}{212--214}
\crow{cosetKernelSum}{\href{https://github.com/lsernau/grh-lean/blob/72fead102f663d5c7a407964a13dc0cf9d977207/grh-compact/src/New/Rows/MemberKernels.lean\#L218-L221}{{\leanfont cosetKernelSum}}}{MemberKernels.lean}{218--221}
\crow{depthindicatorjet}{\href{https://github.com/lsernau/grh-lean/blob/72fead102f663d5c7a407964a13dc0cf9d977207/grh-compact/src/New/Rows/GluedFormTheta.lean\#L33-L143}{{\leanfont Data.\allowbreak{}depth\_indicator\_jet}}}{GluedFormTheta.lean}{33--143}
\crow{DividedLattice}{\href{https://github.com/lsernau/grh-lean/blob/72fead102f663d5c7a407964a13dc0cf9d977207/grh-compact/src/New/Rows/GluedFormTheta.lean\#L147-L151}{{\leanfont DividedLattice}}}{GluedFormTheta.lean}{147--151}
\crow{reducedInverse}{\href{https://github.com/lsernau/grh-lean/blob/72fead102f663d5c7a407964a13dc0cf9d977207/grh-compact/src/New/Rows/GluedFormTheta.lean\#L159-L160}{{\leanfont DividedLattice.\allowbreak{}reducedInverse}}}{GluedFormTheta.lean}{159--160}
\crow{weilScalar}{\href{https://github.com/lsernau/grh-lean/blob/72fead102f663d5c7a407964a13dc0cf9d977207/grh-compact/src/New/Rows/GluedFormTheta.lean\#L164-L165}{{\leanfont DividedLattice.\allowbreak{}weilScalar}}}{GluedFormTheta.lean}{164--165}
\crow{reduceAtconj}{\href{https://github.com/lsernau/grh-lean/blob/72fead102f663d5c7a407964a13dc0cf9d977207/grh-compact/src/New/Rows/GluedFormTheta.lean\#L172-L191}{{\leanfont DividedLattice.\allowbreak{}reduceAt\_conj\_eq\_one}}}{GluedFormTheta.lean}{172--191}
\crow{existsreducedInverse}{\href{https://github.com/lsernau/grh-lean/blob/72fead102f663d5c7a407964a13dc0cf9d977207/grh-compact/src/New/Rows/GluedFormTheta.lean\#L194-L208}{{\leanfont DividedLattice.\allowbreak{}exists\_reducedInverse\_eq}}}{GluedFormTheta.lean}{194--208}
\crow{stackOverGenScalar}{\href{https://github.com/lsernau/grh-lean/blob/72fead102f663d5c7a407964a13dc0cf9d977207/grh-compact/src/New/Rows/GluedFormTheta.lean\#L212-L213}{{\leanfont DividedSplitting.\allowbreak{}stackOverGenScalar}}}{GluedFormTheta.lean}{212--213}
\crow{rowFixingLevelle}{\href{https://github.com/lsernau/grh-lean/blob/72fead102f663d5c7a407964a13dc0cf9d977207/grh-compact/src/New/Rows/GluedFormTheta.lean\#L219-L222}{{\leanfont rowFixingLevel\_le\_levelOf}}}{GluedFormTheta.lean}{219--222}
\crow{congruenceGroup}{\href{https://github.com/lsernau/grh-lean/blob/72fead102f663d5c7a407964a13dc0cf9d977207/grh-compact/src/New/Rows/GluedFormTheta.lean\#L226-L227}{{\leanfont DividedSplitting.\allowbreak{}congruenceGroup}}}{GluedFormTheta.lean}{226--227}
\crow{reduceAtliftMod}{\href{https://github.com/lsernau/grh-lean/blob/72fead102f663d5c7a407964a13dc0cf9d977207/grh-compact/src/New/Rows/GluedFormTheta.lean\#L230-L232}{{\leanfont DividedSplitting.\allowbreak{}reduceAt\_liftMod\_congruenceGroup}}}{GluedFormTheta.lean}{230--232}
\crow{IsGlued}{\href{https://github.com/lsernau/grh-lean/blob/72fead102f663d5c7a407964a13dc0cf9d977207/grh-compact/src/New/Rows/GluedFormTheta.lean\#L236-L242}{{\leanfont DividedSplitting.\allowbreak{}IsGlued}}}{GluedFormTheta.lean}{236--242}
\crow{productForm}{\href{https://github.com/lsernau/grh-lean/blob/72fead102f663d5c7a407964a13dc0cf9d977207/grh-compact/src/New/Rows/GluedFormTheta.lean\#L246-L248}{{\leanfont DividedSplitting.\allowbreak{}productForm}}}{GluedFormTheta.lean}{246--248}
\crow{existscongruenceGroupinv}{\href{https://github.com/lsernau/grh-lean/blob/72fead102f663d5c7a407964a13dc0cf9d977207/grh-compact/src/New/Rows/GluedFormTheta.lean\#L251-L256}{{\leanfont DividedSplitting.\allowbreak{}exists\_congruenceGroup\_inv\_eq\_one\_sub}}}{GluedFormTheta.lean}{251--256}
\crow{existsisGlued}{\href{https://github.com/lsernau/grh-lean/blob/72fead102f663d5c7a407964a13dc0cf9d977207/grh-compact/src/New/Rows/GluedFormTheta.lean\#L259-L294}{{\leanfont DividedSplitting.\allowbreak{}exists\_isGlued}}}{GluedFormTheta.lean}{259--294}
\crow{thetapointFormL}{\href{https://github.com/lsernau/grh-lean/blob/72fead102f663d5c7a407964a13dc0cf9d977207/grh-compact/src/New/Rows/GluedFormTheta.lean\#L308-L336}{{\leanfont DividedSplitting.\allowbreak{}theta\_pointForm}}}{GluedFormTheta.lean}{308--336}
\crow{thetaproductForm}{\href{https://github.com/lsernau/grh-lean/blob/72fead102f663d5c7a407964a13dc0cf9d977207/grh-compact/src/New/Rows/GluedFormTheta.lean\#L343-L366}{{\leanfont DividedSplitting.\allowbreak{}theta\_productForm}}}{GluedFormTheta.lean}{343--366}
\crow{thetastdClassSumproductForm}{\href{https://github.com/lsernau/grh-lean/blob/72fead102f663d5c7a407964a13dc0cf9d977207/grh-compact/src/New/Rows/GluedFormTheta.lean\#L372-L378}{{\leanfont DividedSplitting.\allowbreak{}theta\_stdClassSum\_productForm}}}{GluedFormTheta.lean}{372--378}
\crow{normconjintAlgPart}{\href{https://github.com/lsernau/grh-lean/blob/72fead102f663d5c7a407964a13dc0cf9d977207/grh-compact/src/New/Rows/MonomialCoefficientVector.lean\#L20-L36}{{\leanfont Data.\allowbreak{}Pair.\allowbreak{}norm\_conj\_intAlgPart\_le}}}{MonomialCoefficientVector.lean}{20--36}
\crow{iotagoodProjTranslate}{\href{https://github.com/lsernau/grh-lean/blob/72fead102f663d5c7a407964a13dc0cf9d977207/grh-compact/src/New/Rows/MonomialCoefficientVector.lean\#L38-L41}{{\leanfont Data.\allowbreak{}iota\_goodProjTranslate\_apply}}}{MonomialCoefficientVector.lean}{38--41}
\crow{goodProjTranslateintegral}{\href{https://github.com/lsernau/grh-lean/blob/72fead102f663d5c7a407964a13dc0cf9d977207/grh-compact/src/New/Rows/MonomialCoefficientVector.lean\#L47-L54}{{\leanfont Data.\allowbreak{}goodProjTranslate\_mem\_integralForms}}}{MonomialCoefficientVector.lean}{47--54}
\crow{normSqAtgoodProjTranslate}{\href{https://github.com/lsernau/grh-lean/blob/72fead102f663d5c7a407964a13dc0cf9d977207/grh-compact/src/New/Rows/MonomialCoefficientVector.lean\#L58-L70}{{\leanfont Data.\allowbreak{}normSqAt\_goodProjTranslate\_le}}}{MonomialCoefficientVector.lean}{58--70}
\crow{existsintvector}{\href{https://github.com/lsernau/grh-lean/blob/72fead102f663d5c7a407964a13dc0cf9d977207/grh-compact/src/New/Rows/MonomialCoefficientVector.lean\#L80-L94}{{\leanfont exists\_int\_vector\_harmonicPoly\_eq}}}{MonomialCoefficientVector.lean}{80--94}
\crow{descentsVanishhypSum}{\href{https://github.com/lsernau/grh-lean/blob/72fead102f663d5c7a407964a13dc0cf9d977207/grh-compact/src/New/Rows/MemberOrthogonality.lean\#L35-L53}{{\leanfont descentsVanish\_hypSum}}}{MemberOrthogonality.lean}{35--53}
\crow{stdClassSumperiod}{\href{https://github.com/lsernau/grh-lean/blob/72fead102f663d5c7a407964a13dc0cf9d977207/grh-compact/src/New/Rows/MemberOrthogonality.lean\#L62-L120}{{\leanfont stdClassSum\_period\_eq\_zero}}}{MemberOrthogonality.lean}{62--120}
\crow{stdClassSuminnerproductForm}{\href{https://github.com/lsernau/grh-lean/blob/72fead102f663d5c7a407964a13dc0cf9d977207/grh-compact/src/New/Rows/MemberOrthogonality.lean\#L126-L138}{{\leanfont DividedSplitting.\allowbreak{}stdClassSum\_inner\_productForm\_eq\_zero}}}{MemberOrthogonality.lean}{126--138}
\crow{loss}{\href{https://github.com/lsernau/grh-lean/blob/72fead102f663d5c7a407964a13dc0cf9d977207/grh-compact/src/New/Rows/CartanIntegrality.lean\#L19-L20}{{\leanfont DividedSplitting.\allowbreak{}loss}}}{CartanIntegrality.lean}{19--20}
\crow{normschurinvcols}{\href{https://github.com/lsernau/grh-lean/blob/72fead102f663d5c7a407964a13dc0cf9d977207/grh-compact/src/New/Rows/CartanIntegrality.lean\#L29-L51}{{\leanfont DividedSplitting.\allowbreak{}norm\_schur\_inv\_cols\_le\_weylGauge}}}{CartanIntegrality.lean}{29--51}
\crow{weylGaugecartanProduct}{\href{https://github.com/lsernau/grh-lean/blob/72fead102f663d5c7a407964a13dc0cf9d977207/grh-compact/src/New/Rows/CartanIntegrality.lean\#L56-L69}{{\leanfont DividedSplitting.\allowbreak{}weylGauge\_cartanProduct\_le}}}{CartanIntegrality.lean}{56--69}
\crow{smallHookNumber}{\href{https://github.com/lsernau/grh-lean/blob/72fead102f663d5c7a407964a13dc0cf9d977207/grh-compact/src/New/Rows/ComparisonForms.lean\#L17-L17}{{\leanfont smallHookNumber}}}{ComparisonForms.lean}{17--17}
\crow{SmallFormValues}{\href{https://github.com/lsernau/grh-lean/blob/72fead102f663d5c7a407964a13dc0cf9d977207/grh-compact/src/New/Rows/ComparisonForms.lean\#L21-L22}{{\leanfont SmallFormValues}}}{ComparisonForms.lean}{21--22}
\crow{Tables}{\href{https://github.com/lsernau/grh-lean/blob/72fead102f663d5c7a407964a13dc0cf9d977207/grh-compact/src/New/Rows/ComparisonForms.lean\#L32-L32}{{\leanfont Data.\allowbreak{}Tables}}}{ComparisonForms.lean}{32--32}
\crow{splittingV}{\href{https://github.com/lsernau/grh-lean/blob/72fead102f663d5c7a407964a13dc0cf9d977207/grh-compact/src/New/Rows/ComparisonForms.lean\#L38-L42}{{\leanfont splittingV}}}{ComparisonForms.lean}{38--42}
\crow{classAutrowFixingLevel}{\href{https://github.com/lsernau/grh-lean/blob/72fead102f663d5c7a407964a13dc0cf9d977207/grh-compact/src/New/Rows/ComparisonForms.lean\#L46-L49}{{\leanfont classAut\_rowFixingLevel\_eq\_bot}}}{ComparisonForms.lean}{46--49}
\crow{glueLattice}{\href{https://github.com/lsernau/grh-lean/blob/72fead102f663d5c7a407964a13dc0cf9d977207/grh-compact/src/New/Rows/ComparisonForms.lean\#L53-L54}{{\leanfont glueLattice}}}{ComparisonForms.lean}{53--54}
\crow{lineScale}{\href{https://github.com/lsernau/grh-lean/blob/72fead102f663d5c7a407964a13dc0cf9d977207/grh-compact/src/New/Rows/ComparisonForms.lean\#L59-L59}{{\leanfont lineScale}}}{ComparisonForms.lean}{59--59}
\crow{comparisonForm}{\href{https://github.com/lsernau/grh-lean/blob/72fead102f663d5c7a407964a13dc0cf9d977207/grh-compact/src/New/Rows/ComparisonForms.lean\#L70-L75}{{\leanfont comparisonForm}}}{ComparisonForms.lean}{70--75}
\crow{inneradjoint}{\href{https://github.com/lsernau/grh-lean/blob/72fead102f663d5c7a407964a13dc0cf9d977207/grh-compact/src/New/Rows/ComparisonForms.lean\#L81-L111}{{\leanfont inner\_adjoint\_comparisonForm}}}{ComparisonForms.lean}{81--111}
\crow{comparisonFormintegral}{\href{https://github.com/lsernau/grh-lean/blob/72fead102f663d5c7a407964a13dc0cf9d977207/grh-compact/src/New/Rows/ComparisonForms.lean\#L116-L126}{{\leanfont comparisonForm\_mem\_integralForms}}}{ComparisonForms.lean}{116--126}
\crow{normSqAtcomparisonForm}{\href{https://github.com/lsernau/grh-lean/blob/72fead102f663d5c7a407964a13dc0cf9d977207/grh-compact/src/New/Rows/ComparisonForms.lean\#L131-L163}{{\leanfont normSqAt\_comparisonForm\_le}}}{ComparisonForms.lean}{131--163}
\crow{tableThetaadjoint}{\href{https://github.com/lsernau/grh-lean/blob/72fead102f663d5c7a407964a13dc0cf9d977207/grh-compact/src/New/Rows/ComparisonForms.lean\#L170-L192}{{\leanfont tableTheta\_adjoint\_eq\_zero}}}{ComparisonForms.lean}{170--192}
\crow{existscomparisonforms}{\href{https://github.com/lsernau/grh-lean/blob/72fead102f663d5c7a407964a13dc0cf9d977207/grh-compact/src/New/Rows/ComparisonForms.lean\#L200-L227}{{\leanfont exists\_comparison\_forms}}}{ComparisonForms.lean}{200--227}}
% the notations of the defined functions: one slot for each explicit argument of the Lean declaration
\newcommand{\pointForm}[4]{\operatorname{pt}_{#1,#2,#3}(#4)}
\newcommand{\kernelVector}[8]{\operatorname{kv}_{#1,#2;#3}[#4;#5](#6,#7)_{#8}}
\newcommand{\rowFixingLevel}[3]{\Gamma(#1;#2,#3)}
\newcommand{\DividedLattice}[1]{\operatorname{Lat}(#1)}
\newcommand{\reducedInverse}[1]{G_{#1}}
\newcommand{\weilScalar}[1]{\beta_{#1}}
\newcommand{\stackOverGenScalar}[2]{[#2;\bar\varpi]_{#1}}
\newcommand{\congruenceGroup}[1]{\mathsf{C}(#1)}
\newcommand{\IsGlued}[3]{\operatorname{Glued}_{#1}(#2,#3)}
\newcommand{\productForm}[4]{\operatorname{pr}_{#1,#2}(#3;#4)}
\newcommand{\loss}[1]{\operatorname{loss}(#1)}
\newcommand{\smallHookNumber}[1]{\mathsf{h}_{D}(#1)}
\newcommand{\SmallFormValues}[1]{\mathcal{W}_{D}(#1)}
\newcommand{\Tables}[1]{\operatorname{Tab}_{D}(#1)}
\newcommand{\splittingV}[2]{\mathcal{V}_{D}(#1,#2)}
\newcommand{\glueLattice}[3]{L_{D}(#1,#2,#3)}
\newcommand{\lineScale}[1]{\beta^{\mathrm{line}}_{D}(#1)}
\newcommand{\comparisonForm}[3]{\operatorname{gl}_{D,#1,#2}(#3)}
\newcommand{\UnitMembers}{219}
\newcommand{\UnitInNote}{57}
\newcommand{\UnitProofs}{85}
\newcommand{\UnitSteps}{634}
\newcommand{\UnitAllSteps}{680}


\linespread{1.08}                       % 11 pt on 14.7 pt
\setlength{\parindent}{0pt}
\setlength{\parskip}{5.5pt plus 1.5pt minus 0.5pt}
\setlength{\abovedisplayskip}{7.5pt plus 2pt minus 3pt}
\setlength{\belowdisplayskip}{7.5pt plus 2pt minus 3pt}
\setlength{\abovedisplayshortskip}{3pt plus 2pt}
\setlength{\belowdisplayshortskip}{6pt plus 2pt minus 2pt}
\emergencystretch=1.5em
\clubpenalty=10000 \widowpenalty=10000 \displaywidowpenalty=10000
\relpenalty=7000 \binoppenalty=9000
\raggedbottom
\pagestyle{fancy}\fancyhf{}\renewcommand{\headrulewidth}{0pt}
\fancyfoot[C]{\sffamily\fontsize{7.6}{9}\selectfont\color{inkfaint}\thepage\,/\,\pageref*{LastPage}}
\setlist{topsep=2pt,partopsep=0pt,parsep=0pt,itemsep=4pt,leftmargin=2.1em,labelsep=0.55em}

\newcommand{\Q}{\mathbb{Q}}
\newcommand{\Z}{\mathbb{Z}}
\newcommand{\F}{\mathbb{F}}
\newcommand{\len}{\operatorname{len}}
\newcommand{\cok}{\operatorname{cok}}
\newcommand{\End}{\operatorname{End}}
\newcommand{\Hom}{\operatorname{Hom}}
\newcommand{\Gal}{\operatorname{Gal}}
\newcommand{\diag}{\operatorname{diag}}
\newcommand{\kw}[1]{\emph{#1}}
\newcommand{\m}[1]{\mbox{$#1$}}        % a short formula that must not be broken across lines
\newcommand{\head}[1]{\par\addvspace{17pt plus 4pt minus 2pt}\Needspace{6\baselineskip}%
  \noindent{\fontsize{13.2}{16}\selectfont\bfseries\color{ink}#1}\par\nopagebreak\vspace{-1pt}}
\newcommand{\sect}[2]{\refstepcounter{section}\label{#2}\par\addvspace{19pt plus 4pt minus 2pt}\Needspace{7\baselineskip}%
  \noindent\llap{{\fontsize{13.2}{16}\selectfont\bfseries\color{inkfaint}\thesection}\hspace{9pt}}%
  {\fontsize{13.2}{16}\selectfont\bfseries\color{ink}#1}\par\nopagebreak\vspace{-1pt}}
\newcommand{\subhead}[1]{\par\addvspace{9pt plus 2pt}\Needspace{4\baselineskip}\noindent\textbf{#1}\enspace\ignorespaces}
% where a statement lives in the Lean archive
\newcommand{\leantag}{{\sffamily\bfseries\fontsize{6.3}{8}\selectfont\color{leancol}\addfontfeatures{LetterSpace=7}LEAN}}
\newcommand{\sep}{\unskip\nobreak\hspace{5pt}{\color{inkfaint}\textperiodcentered}\hspace{5pt plus 2pt}\ignorespaces}
\newcommand{\leanline}[1]{\par\nopagebreak\vspace{-2.5pt}{\fontsize{8.9}{12.4}\selectfont\color{inksoft}\raggedright
  \noindent\llap{\leantag\hspace{9pt}}#1\par}\vspace{1pt}}
% a displayed statement in a tinted box: #1 colour, #2 text
\newcommand{\statement}[2]{\par\addvspace{7pt}\begingroup\setlength{\fboxsep}{8.5pt}\setlength{\fboxrule}{0.6pt}%
  \noindent\fcolorbox{#1!30!white}{#1!3!white}{\parbox{\dimexpr\linewidth-2\fboxsep-2\fboxrule\relax}{#2}}%
  \endgroup\par\addvspace{7pt}}
\newcommand{\nt}[2]{\par\noindent\hangindent=1.1em\hangafter=1 #1\enspace{\color{inksoft}#2}\par\vspace{1pt}}
\newcommand{\legdot}[1]{\tikz[baseline=-0.62ex]\fill[#1] (0,0) circle (2.5pt);\hspace{3pt}}
\newcommand{\crow}[4]{\csname CS@#1\endcsname & #2 & \textcolor{inksoft}{#3} & \textcolor{inksoft}{#4}\\}
\newcommand{\leanquoteline}[2]{\par\noindent\hangindent=\dimexpr#1\fontdimen2\font+2.2em\relax\hangafter=1
  \hspace*{#1\fontdimen2\font}#2}
\newcommand{\quoted}[1]{\par\vspace{3pt}{\ttfamily\fontsize{7.5}{10.4}\selectfont\color{leancol}\raggedright
  \setlength{\parskip}{0pt}\leanquote{#1}\par}\vspace{4pt}}
\newcommand{\bx}[3]{{\sffamily\bfseries\fontsize{7.3}{9.2}\selectfont\color{#1!82!black}#2}\\[1.5pt]#3}
\newcommand{\cs}[2]{\expandafter\newcommand\csname CS@#1\endcsname{#2}}

% theorems
\newtheoremstyle{plainx}{7pt plus 2pt}{7pt plus 2pt}{\itshape}{}{\bfseries}{.}{0.5em}{}
\newtheoremstyle{uprightx}{7pt plus 2pt}{7pt plus 2pt}{}{}{\bfseries}{.}{0.5em}{}
\theoremstyle{plainx}
\newtheorem{theorem}{Theorem}[section]
\newtheorem{lemma}[theorem]{Lemma}
\newtheorem{proposition}[theorem]{Proposition}
\newtheorem{corollary}[theorem]{Corollary}
\theoremstyle{uprightx}
\newtheorem{definition}[theorem]{Definition}
\newtheorem{setting}[theorem]{Setting}          % hypotheses shared by several statements; each of them cites it by number
\newtheorem{remark}[theorem]{Remark}
\newtheorem{example}[theorem]{Example}
\renewcommand{\qedsymbol}{{\setlength{\fboxsep}{0pt}\setlength{\fboxrule}{0.5pt}\fbox{\rule{0pt}{0.46em}\hspace{0.46em}}}}
% what holds throughout the note: the ambient objects only (implicit and instance arguments), never a hypothesis
\newenvironment{conventions}{\par\addvspace{5pt}\noindent\textbf{Conventions.}\enspace\ignorespaces}{\par\addvspace{3pt}}
\makeatletter
\renewenvironment{proof}[1][Proof]{\par\pushQED{\qed}\normalfont\topsep0pt\partopsep0pt
  \trivlist\item[\hskip\labelsep\itshape #1.]\ignorespaces}{\popQED\endtrivlist\@endpefalse}
\makeatother

% where each declaration of the concordance is treated: \cs{key}{place}, the place naming a \label or a lettered item
\cs{sumdiagramWeight}{Lemma~\ref{l:small}}
\cs{existsalgEquiv}{Lemma~\ref{l:small}}
\cs{normjetle}{Proposition~\ref{p:jet}}
\cs{boxesle}{Lemma~\ref{l:small}}
\cs{existsisHCParamofle}{Lemma~\ref{l:small}}
\cs{existsisHCParamofhasGap}{Lemma~\ref{l:small}}
\cs{modKernelembFin}{Lemma~\ref{l:small}}
\cs{summemintegralAbove}{Lemma~\ref{l:small}}
\cs{existsmodKernelbound}{Proposition~\ref{p:crude}}
\cs{existsreduce}{Lemma~\ref{l:reduce}}
\cs{pointForm}{Definition~\ref{d:point}}
\cs{innerpointForm}{Lemma~\ref{l:inner}}
\cs{thetapointForm}{Lemma~\ref{l:thetapoint}}
\cs{kernelVector}{Definition~\ref{d:point}}
\cs{rowFixingLevel}{Definition~\ref{d:point}}
\cs{cosetKernelSum}{Definition~\ref{d:point}}
\cs{depthindicatorjet}{Proposition~\ref{p:depth}}
\cs{DividedLattice}{Definition~\ref{d:glue}}
\cs{reducedInverse}{Definition~\ref{d:glue}}
\cs{weilScalar}{Definition~\ref{d:glue}}
\cs{reduceAtconj}{Lemma~\ref{l:conj}}
\cs{existsreducedInverse}{Lemma~\ref{l:latexists}}
\cs{stackOverGenScalar}{Definition~\ref{d:glue}}
\cs{rowFixingLevelle}{Lemma~\ref{l:rowfix}}
\cs{congruenceGroup}{Definition~\ref{d:glue}}
\cs{reduceAtliftMod}{Remark~\ref{r:liftC}}
\cs{IsGlued}{Definition~\ref{d:glue}}
\cs{productForm}{Definition~\ref{d:glue}}
\cs{existscongruenceGroupinv}{Lemma~\ref{l:rowfix}}
\cs{existsisGlued}{Proposition~\ref{p:glue}}
\cs{thetapointFormL}{Lemma~\ref{l:thetaL}}
\cs{thetaproductForm}{Proposition~\ref{p:thetaprod}}
\cs{thetastdClassSumproductForm}{Proposition~\ref{p:thetaprod}}
\cs{normconjintAlgPart}{Lemma~\ref{l:cleared}}
\cs{iotagoodProjTranslate}{Lemma~\ref{l:cleared}}
\cs{goodProjTranslateintegral}{Lemma~\ref{l:cleared}}
\cs{normSqAtgoodProjTranslate}{Lemma~\ref{l:cleared}}
\cs{existsintvector}{Lemma~\ref{l:cleared}}
\cs{descentsVanishhypSum}{Lemma~\ref{l:hyp}}
\cs{stdClassSumperiod}{Proposition~\ref{p:period}}
\cs{stdClassSuminnerproductForm}{Corollary~\ref{c:orth}}
\cs{loss}{Definition~\ref{d:glue}}
\cs{normschurinvcols}{Lemma~\ref{l:cartan}}
\cs{weylGaugecartanProduct}{Lemma~\ref{l:cartan}}
\cs{smallHookNumber}{Definition~\ref{d:comp}}
\cs{SmallFormValues}{Definition~\ref{d:comp}}
\cs{Tables}{Definition~\ref{d:comp}}
\cs{splittingV}{Definition~\ref{d:comp}}
\cs{classAutrowFixingLevel}{Remark~\ref{r:liftC}}
\cs{glueLattice}{Definition~\ref{d:comp}}
\cs{lineScale}{Definition~\ref{d:comp}}
\cs{comparisonForm}{Definition~\ref{d:comp}}
\cs{inneradjoint}{Lemma~\ref{l:adj}}
\cs{comparisonFormintegral}{Lemma~\ref{l:intsize}}
\cs{normSqAtcomparisonForm}{Lemma~\ref{l:intsize}}
\cs{tableThetaadjoint}{Lemma~\ref{l:tt}}
\cs{existscomparisonforms}{Theorem A}
\newcommand{\CC}{\mathbb{C}}
\newcommand{\ZZ}{\mathbb{Z}}
\newcommand{\NN}{\mathbb{N}}
\newcommand{\AAA}{\mathbb{A}}
\newcommand{\ix}[1]{\mathsf{#1}}
\newcommand{\OEr}{\mathcal{O}_E}
\newcommand{\Ohat}{\widehat{\mathcal{O}}_E}
\newcommand{\UU}{\mathrm{U}}
\newcommand{\GL}{\mathrm{GL}}
\newcommand{\KK}[1]{\mathbf{K}_{#1}}
\newcommand{\ip}[2]{\langle #1,#2\rangle}
\newcommand{\clsum}[2]{\sideset{}{^{\mathrm{cl}}}\sum_{#1\bmod #2}}
\newcommand{\tk}[7]{\theta^{#1,#2}_{#3,#4;#5}[#6,#7]}
\newcommand{\AutC}{\operatorname{Aut}(\CC)}
\newcommand{\transp}{^{\mathsf{T}}}
\newcommand{\herm}{^{\mathsf{H}}}
\newcommand{\inputbox}[1]{\statement{inksoft}{#1}}
% objects of the sister notes, written by hand with all their arguments
\newcommand{\FormVx}[1]{\mathcal{F}(#1)}
\newcommand{\dsc}[1]{\operatorname{lift}_{#1}}
\newcommand{\gau}[5]{\operatorname{gauge}^{#1}_{#2}(#3;#4,#5)}
\newcommand{\nsq}[2]{\|#1\|^{2}_{#2}}
\newcommand{\cks}[5]{\operatorname{cks}_{#1,#2}(#3,#4;#5)}
\newcommand{\Mgood}[1]{\mathcal{M}^{\circ}_D(W_{#1}(D))}
\newcommand{\alift}{\operatorname{lift}^{*}}
\newcommand{\lmult}[4]{\operatorname{lm}_D(#1,#2,#3,#4)}
\newcommand{\fnat}[2]{F^{\natural}_{D}(#1,#2)}
\newcommand{\ttheta}[2]{\operatorname{tt}_D(#1,#2)}
\newcommand{\HV}[2]{\operatorname{HV}_D(#1,#2)}
\newcommand{\Mem}[1]{\operatorname{Mem}_D(#1)}
\newcommand{\Bix}[1]{\mathsf{B}_D(#1)}
\newcommand{\Six}{\mathsf{S}_D}
\newcommand{\Cix}{\mathsf{C}_D}
\newcommand{\Scl}{\operatorname{Cl}^{\circ}(D)}

\begin{document}
\color{ink}

{\raggedright
{\sffamily\bfseries\fontsize{7.2}{9}\selectfont\color{leancol}\addfontfeatures{LetterSpace=8}READBACK NOTE · CHAPTER VII · UNIT 143 · PART 3 OF 4\par}
\vspace{5pt}
{\fontsize{21.5}{25.5}\selectfont\bfseries Glued lattices, product forms and the integral comparison forms\par}
\vspace{5pt}
{\fontsize{12.4}{16}\selectfont\itshape\color{inksoft} One unit of the Lean artifact \ArtifactName, read back into
mathematics\par}}
\vspace{8pt}
{\color{rulegrey}\hrule height 0.6pt}
\vspace{7pt}
\thispagestyle{empty}

\textbf{Summary.} To bound the volume of a single theta lift at a twist against the volume at the trivial twist, the
proof needs forms on the definite side that can be compared for all twists at once. This note constructs them: for
each lift index~$j$ and each twist~$\vartheta$ a linear map from a fixed space of \emph{tables} to the algebraic
forms of the auxiliary pair of \m{(\vartheta,j)}, such that the adjoint lift carries all these maps to one and the
same surjection onto the good forms, up to a non-zero number \m{\gamma(\vartheta)}; the forms are integral away from
the conductor of the twist, small at the trivial twist, and orthogonal to what the congruence bound of Chapter~VIII
requires (Theorem~A). The forms are sums of delta functions at the points of \emph{glued lattices}, with values
products of two coordinate vectors, and their descents are averages of products of two theta kernels. Along the way
the note proves crude integrality and size bounds for theta kernels of definite pairs, the vanishing of theta periods
for forms that are not lifts, and that one explicit jet of a theta function is no more divisible than predicted.

\statement{leancol}{\fontsize{9.6}{13.2}\selectfont\textbf{What this note is.} The artifact \ArtifactName{} is a
body of Lean code: the compact edition of a formal proof, due to L.~Alpöge with the AI system Claude, of the Riemann
hypothesis and of its generalisation to Dirichlet $L$-functions. That headline is a claim. The authors report that
every file compiles with Lean v4.35.0 against a pinned version of Mathlib and that Lean's kernel accepted the
proofs; an independent build is recorded in the handout's record of checks when it reports. No mathematician has
read the argument through, and nothing here says that the Riemann hypothesis is proved. This note is the third of
four on one unit of the artifact, the folder of the rows; it treats \UnitInNote{} of the unit's \UnitMembers{}
declarations, those of six files. Every declaration is stated. The proofs are given at the level at which they were
read: complete for the algebra of glued lattices and for the assembly of Theorem~A, and as an account of which
inputs combine for the analytic and $p$-adic lemmas, whose inputs belong to other units and to the background
library and are named and not proved here. Every statement was compared with its Lean text, argument by argument.
The Lean text read is \PinText; Appendix~A is the concordance of \ConcDecls{} declarations.}

\sect{The statement}{s:intro}
\statement{arith}{\phantomsection\label{t:A}\textbf{Theorem A} (comparison forms).\enspace
\emph{In Settings \ref{set:bg} and~\ref{set:data}, let $P$ be parameters, $D$ admissible data for~$P$ and $j$ a lift
index. Then there are a function~$\gamma$ from the twists of~$D$ to~$\CC$, a linear map
\m{A\colon\Tables{j}\to\Mgood{j}} and, for every twist~$\vartheta$ of~$D$, a linear map
\m{\operatorname{gl}_\vartheta\colon\Tables{j}\to\FormVx{O_D(\vartheta,j)}} with the following properties, in which
\m{u_\vartheta} denotes the restriction of the adjoint lift \m{\alift_D(\vartheta,j)} to \m{\Mgood{j}} and
\m{u_\vartheta^{*}} its adjoint.
\begin{enumerate}[label=\textup{(\roman*)},leftmargin=2.3em,itemsep=2pt]
\item For every twist~$\vartheta$: \m{\gamma(\vartheta)\ne0} and
\m{|\gamma(\vartheta)/\gamma(1)|^2=\HV{j}{\vartheta}/\HV{j}{1}}.
\item For every twist~$\vartheta$: \m{u_\vartheta^{*}\circ\operatorname{gl}_\vartheta=\gamma(\vartheta)\cdot A}.
\item $A$ is surjective.
\item For every twist~$\vartheta$ and every \m{y\in\Mgood{j}}, the form \m{\operatorname{gl}_\vartheta(A^{*}y)} is tested
against the orthogonal complement of the image of~$u_\vartheta$.
\item For every twist~$\vartheta$, every table~$\zeta$, every prime~$\ell$ prime to the norms of the primes in
\m{Q(\vartheta)} and every embedding \m{\varphi\colon\CC\to\CC_\ell}: if, for every index~$a$ of the table, the
conjugate vector \m{\overline{\zeta_a}} has gauge at most~$1$ at the representative of its class for the embedding
\m{\varphi\circ\mathrm{conj}}, then \m{\operatorname{gl}_\vartheta(\zeta)} is integral at~$\varphi$.
\item For every table~$\zeta$ and every \m{\sigma\in\AutC}:
\[\smallHookNumber{j}\cdot\nsq{\operatorname{gl}_1(\zeta)}{\sigma}\le\exp\bigl(2\,r\,B_0(D)\bigr)\cdot\sum_a\|\sigma(\zeta_a)\|^2 .\]
\end{enumerate}}}
\leanline{\lean{existscomparisonforms}}

Here \m{\Tables{j}} is the space of tables (Definition~\ref{d:comp}), \m{\HV{j}{\vartheta}} is a product of Hecke
$L$-values at twisted characters (Setting~\ref{set:data}), \m{B_0(D)} is the budget for congruences, and ``tested'',
``gauge'' and ``integral'' are recalled in Setting~\ref{set:bg}\,(d). In (i) the twist enters only through
\m{\gamma(\vartheta)}: this is what ``the comparison forms see the twist through one number'' means. The note
\emph{The height of a ratio of lift volumes} deduces from (i) to~(vi) the height bound for a ratio of lift volumes.

\sect{The setting}{s:setting}
\begin{conventions}\phantomsection\label{c:conventions}
Finite index sets are written in sans-serif with the letters of the formal text, as in the note \emph{Auxiliary
dual pairs, determinant polynomials and the carrier}; \m{|\ix n|} is the number of elements. The ambient objects of
the declarations are: a field~$E$ that is a CM field with a CM type and, where theta kernels occur, with Weil data;
the finite index sets \m{\ix r,\ix n,\ix N,\ix m,\ix a,\ix l,\ix L,\ix c} of each statement; an auxiliary pair~$O$
for \m{(E,\ix r,\ix n)} and an embedding \m{z\colon\ix r\hookrightarrow\ix m}, where a divided splitting occurs; in
the statements about data the parameters~$P$ and the data~$D$; a prime~$p$ with an embedding
\m{\iota_p\colon\CC\to\CC_p}, where gauges occur. In statements whose formal text introduces further implicit
objects (characters, moduli, vectors, points), these are introduced in the statement itself.
\end{conventions}

\begin{setting}[background and auxiliary pairs; from the background library and the sister note]\label{set:bg}
\emph{(a)} $E$ is a CM field with a CM type and Weil data; $\upsilon$, moduli, characters with weight and tame
part, hermitian spaces, the groups \m{\UU(S)(\AAA)} and \m{G_\infty(s)}, the stabiliser \m{\KK{\ix n}} of the standard
lattice of \m{I_{\ix n}}, levels, class sums \m{\clsum{h}{\Gamma}}, the theta kernel
\m{\tk{S}{T}{\alpha}{\beta}{\mathfrak{m}}{\Psi}{P}(g,h)}, weights~$\kappa$, standard fillings, the polynomial
\m{\operatorname{pol}_\kappa(\xi)}, the Gram norm \m{\|\xi\|_\kappa}, the Fock pairing and the Schur matrix
\m{S_\kappa(z)} are as in the background Setting of the sister note. \m{\ip{\xi}{\xi'}_\kappa=\xi\herm G_\kappa\xi'} is the Gram form.

\emph{(b)} An auxiliary pair~$O$ for \m{(E,\ix r,\ix n)}, with its products of chosen and of twist primes
\m{\mathfrak{M}_O}, \m{\mathfrak{Q}_O}, its space~$W_O$ with signs~$s_O$, its embedding~$e_O$, its numbering
\m{\operatorname{ord}_O}, its weights \m{\kappa^{+}_O}, \m{\kappa^{-}_O}, \m{\kappa_O}, its determinant exponent
\m{\epsilon_O}, its character~$\mu_O$ with tame part~$\psi_O$, its characters \m{\mathbf{a}_O}, \m{\mathbf{b}_O}, its
level~$\Gamma_O$ with character~$\chi_O$, its forms \m{\FormVx{O}}, its harmonic polynomial \m{P^{O}(\xi)}, its descent
\m{\dsc{O}(\Xi)(g)}, its Hecke test operators \m{H_{O,w}(X)} and test values \m{z_O(w)}, the generator \m{\varpi_O} of
\m{\mathfrak{M}_O}, a divided splitting~$\mathcal{V}$ in \m{\operatorname{Div}(\ix N,\ix m,\ix a;O,z)} with its bijection
\m{\varepsilon\colon\ix N\sqcup\ix m\to\ix n}, its fixed embedding \m{\phi_{\mathcal{V}}\colon\ix r\hookrightarrow\ix N} and
its divided embedding \m{\ix a\hookrightarrow\ix N}, the class \m{\bar\varpi_{\mathcal{V}}} of the generator in
\m{\OEr/\mathfrak{M}_O^{2}}, the reduction \m{\operatorname{red}_0} from \m{\OEr/\mathfrak{M}_O^{2}\mathfrak{Q}_O} to
\m{\OEr/\mathfrak{M}_O^{2}}, the stabiliser \m{\operatorname{Stab}(\varphi,w,e)} and the column datum
\m{\operatorname{cd}_{\chi,\varphi,w,e}(\sigma)}, and the inclusion matrices~$I_e$ are those of the sister note.
For a weight, the polynomial of a coordinate vector \emph{in the positive letters} (or \emph{negative letters}) of~$O$
is \m{\operatorname{pol}_\kappa} with each variable \m{(\tau,i,h)} replaced by the sum of the \m{x_{\tau,i,c}} over the
\m{c\in\ix r} with sign $+1$ (or~$-1$) and \m{\operatorname{ord}_O(c)=h}; it is written
\m{\operatorname{pol}^{\pm}_{O,\kappa}(\xi)}.

\emph{(c)} For level places with places~$w$, \m{\tilde A} denotes the chosen lift to \m{\KK{\ix n}} of an invertible
matrix~$A$ over a residue ring; it has component~$1$ outside the primes of the modulus and their conjugates.

\emph{(d) Forms, gauges, tests.} An algebraic form of weight \m{(\kappa,\epsilon)} on~$I_{\ix n}$ has values
coordinate vectors of weight~$\kappa$. For \m{\sigma\in\AutC} the number \m{\nsq{\phi}{\sigma}} is the class sum of
the squared Gram norms of the vectors \m{\sigma(\phi(g))}, and \m{\ip{\phi}{\phi'}} is the class sum of the Gram forms
of the values. For a prime~$p$, an embedding \m{\iota_p\colon\CC\to\CC_p}, a finite adelic point~$g$ and a
coordinate vector~$\xi$, the \kw{gauge} \m{\gau{\iota_p}{\kappa}{\epsilon}{g}{\xi}} is the largest $p$-adic absolute value of
a coordinate of \m{g^{-1}\cdot\iota_p(\xi)}, the action being that of weight \m{(\kappa,\epsilon)}; a form is \kw{integral
at~$\iota_p$} if each value has gauge at most~$1$ at its point. The \kw{delta form} at a point~$x$ with value~$v$ is the
form supported on the class of~$x$ that the level and the rational points generate from~$v$. For an idempotent~$e_0$
of the group algebra of the finite group of the chosen places and a subspace~$\mathcal{Z}$ of \m{\FormVx{O}}, a form is
\kw{tested against}~$\mathcal{Z}$ (for~$e_0$) if it is orthogonal to every form of~$\mathcal{Z}$ that lies in the image
of~$e_0$ and is killed by \m{H_{O,w}(z_O(w))} for some finite place \m{w\ne\bar w} with $w$, $\bar w$ not among the
places of~$O$.
\end{setting}

\begin{setting}[the data; from other units and from the sister note on heights]\label{set:data}
$P$ are parameters and $D$ admissible data for~$P$, with field \m{E(D)}, number of lifts~$m$, ranks $k(D)$ and~$r$,
step \m{t(D)=5r}, scales $x$, $x_1$, $b$, second level \m{g(D)=g_0g_1}, the integer \m{\bar\pi_g} of \m{E(D)} with
\m{\bar\pi_g\overline{\bar\pi_g}=g(D)}, whose prime divisors are the chosen primes, the finite group \m{\mathsf{G}(D)} and
its good idempotent \m{e_0(D)}, and \m{B_0(D)=\operatorname{budget}(D)/(m\,k(D))}. A lift index is an integer~$j$ with
\m{1\le j\le m}. A \kw{twist} of~$D$ is a tame character whose tame part is non-trivial at every prime of its
set~\m{Q(\vartheta)}; the trivial character~$1$ is one. For a twist~$\vartheta$ and a lift index~$j$:
\m{O_D(\vartheta,j)} is the auxiliary pair of the $j$-th single lift, for \m{(E(D),\ix r,\ix n)} with \m{|\ix r|=r} and
\m{|\ix n|=r+t(D)}, with space~$W_j(D)$, chosen primes those of~$D$ and twist primes \m{Q(\vartheta)}; \m{\Mgood{j}} is
the space of good forms on~$W_j(D)$; \m{\alift_D(\vartheta,j)} is the adjoint lift, whose adjoint pairing is
\m{\ip{\alift_D(\vartheta,j)y}{\Xi}=\ip{y}{\dsc{O_D(\vartheta,j)}(\Xi)}} (Petersson product on the right).

From the note \emph{The height of a ratio of lift volumes}: the index sets \m{\Bix{j}} (of size \m{r+9m+1}),
\m{\Six} and \m{\Cix} (of size~$3m$), with \m{\ix n=\Bix{j}\sqcup\Six} up to a fixed bijection; the finite set
\m{\Scl} of classes of \m{\UU(I_{\Six})} at the principal level \m{\bar\pi_g^{2}}, with representatives
\m{h_{\mathfrak{c}}}; the finite set \m{\Mem{j}} of \kw{choices}, a choice~$i$ having a first component~$i_1$, a
\m{\Cix\times\Bix{j}} matrix over \m{\OEr/\bar\pi_g^{2}}; the character \m{\mathbf{b}''_D(\vartheta,j)}; the \kw{line
multiplier} \m{\lmult{\vartheta}{j}{i}{h}}, a coordinate vector of weight \m{\kappa^{+}} with rows \m{\Bix{j}}, for a
choice~$i$ and \m{h\in\UU(I_{\Bix{j}})(\AAA_f)}; the function \m{\fnat{j}{i}} on \m{\UU(W_j(D))(\AAA)}; for a
function~$\zeta$ from \m{\Scl} to coordinate vectors of weight~$\kappa^{-}$ with rows~$\Six$, the function
\[\ttheta{j}{\zeta}(g)=\sum_{\mathfrak{c}\in\Scl}\overline{\tk{W_j(D)}{I_{\Six}}{\mathbf{a}}{\upsilon^{m/2-j}}{\bar\pi_g^{2}}{\mathbf{1}_{I}}{\operatorname{pol}^{-}(\overline{\zeta(\mathfrak{c})})}(g,h_{\mathfrak{c}})}\]
on \m{\UU(W_j(D))(\AAA)}, with \m{\mathbf{a}} the first character of \m{O_D(1,j)}, \m{\mathbf{1}_I} the indicator of the
inclusion matrix of \m{\ix r\hookrightarrow\Six} and the polynomial in the negative letters; and the positive real
number \m{\HV{j}{\vartheta}}, a product of products of Hecke $L$-values at three characters twisted
by~$\vartheta^{\pm1}$. Here \m{\kappa^{\pm}} are the weights of \m{O_D(1,j)}, which do not depend on the twist.
\end{setting}

\subsection*{Definitions of this note}
\begin{definition}[levels, kernel vectors and point forms]\label{d:point}
In Setting~\ref{set:bg}:
\begin{enumerate}[label=\textup{(\roman*)},leftmargin=2.3em,itemsep=2pt]
\item For level places~$S$, finite index sets \m{\ix r}, \m{\ix m}, an index set~\m{\ix c}, a \m{\ix c\times\ix m}
matrix~$w$ over \m{\OEr/\mathfrak{M}_S^{2}} and an embedding \m{e\colon\ix r\hookrightarrow\ix m}, the level
\m{\rowFixingLevel{S}{w}{e}} consists of the elements of \m{\KK{\ix m}} whose reduction modulo
\m{\mathfrak{M}_S^{2}\mathfrak{Q}_S} lies in \m{\operatorname{Stab}(\operatorname{red}_0,w,e)}: its columns outside \m{e(\ix r)}
are those of~$1$, and its reduction~$A$ modulo \m{\mathfrak{M}_S^{2}} satisfies \m{wA=w}.
\item For an auxiliary pair~$O'$, a finite index set~$\ix m$, a character~$b$, a modulus~$\mathfrak{m}$, a
function~$\Psi$ on matrices modulo~$\mathfrak{m}$, a weight~$\kappa$, points $g$, $h$ and a standard filling~$i$ of
weight~$\kappa$ with rows~$\ix m$, the \kw{kernel vector} has the coordinates
\[\kernelVector{O'}{b}{\mathfrak{m}}{\Psi}{\kappa}{g}{h}{i}=\overline{\tk{W_{O'}}{I_{\ix m}}{\mathbf{a}_{O'}}{b}{\mathfrak{m}}{\Psi}{b_i^{-}}(g,h)},\]
where \m{b_i^{-}} is the basis polynomial~$b_i$ written in the negative letters of~$O'$.
\item For an auxiliary pair~$O'$, a finite index set~$\ix m$, a character~$b$, an index set~$\ix c$, a
\m{\ix c\times\ix m} matrix~$w$ over \m{\OEr/\mathfrak{M}_{O'}^{2}}, an embedding~$e$ into~$\ix m$, a function~$y$ from
\m{\UU(I_{\ix m})(\AAA_f)} to coordinate vectors of weight~$\kappa^{+}_{O'}$ and a point~$g$ of
\m{\UU(W_{O'})(\AAA)}, the \kw{coset kernel sum} is
\[\cks{O'}{b}{w}{e}{y}(g)=\clsum{h}{\rowFixingLevel{O'}{w}{e}}\tk{W_{O'}}{I_{\ix m}}{\mathbf{a}_{O'}}{b}{\mathfrak{M}^{2}\mathfrak{Q}}{\operatorname{cd}_{\chi,\operatorname{red}_0,w,e}}{\operatorname{pol}^{+}(y(h))}(g,h),\]
where \m{\mathfrak{M}^{2}\mathfrak{Q}} is the modulus of~$O'$, $\chi$ the character of its residue ring given by the
tame part~$\psi_{O'}$, and the polynomial of \m{y(h)} is in the positive letters of~$O'$. The formal text lists a
further point~$h$ among the arguments; see §\ref{s:notes}.
\item For an auxiliary pair~$O$ for \m{(E,\ix r,\ix n)}, a point \m{h\in\UU(I_{\ix n})(\AAA_f)}, a family \m{Y=(Y_w)}
of \m{\ix n\times\ix r} matrices with linearly independent columns over the residue fields of the chosen places, an
embedding \m{e\colon\ix r\hookrightarrow\ix n} and a coordinate vector~$v$ of weight~$\kappa_O$, the \kw{point form}
\m{\pointForm{h}{Y}{e}{v}\in\FormVx{O}} is the delta form with value~$v$ at the point \m{h\,k_{Y,e}}, where
\m{k_{Y,e}\in\KK{\ix n}} is the chosen element of Lemma~\ref{l:reduce}. It is linear in~$v$.
\end{enumerate}
\end{definition}
\leanline{\lean{rowFixingLevel}\sep \lean{kernelVector}\sep \lean{cosetKernelSum}\sep \lean{pointForm}}

\begin{definition}[divided lattices, glued lattices, product forms]\label{d:glue}
In Setting~\ref{set:bg}, let $O$ be an auxiliary pair for \m{(E,\ix r,\ix n)}.
\begin{enumerate}[label=\textup{(\roman*)},leftmargin=2.3em,itemsep=2pt]
\item A \kw{divided lattice}, an element~$L$ of \m{\DividedLattice{O}}, is a point \m{\gamma_L\in\UU(I_{\ix n})(\AAA_f)}
with component~$1$ outside the chosen places of~$O$ and their conjugates such that every entry of
\m{\varpi_O\,\gamma_L^{-1}} is integral.
\item Its \kw{reduced inverse} \m{\reducedInverse{L}} is the \m{\ix n\times\ix n} matrix over
\m{\OEr/\mathfrak{M}_O^{2}} of the reductions of the entries of \m{\varpi_O\gamma_L^{-1}}, and its \kw{Weil scalar} is
\m{\weilScalar{L}=\chi(\delta)/|\delta|^{|\ix r|/2}}, where $\delta$ is the idele equal to \m{\det\gamma_L} at the chosen
places and to~$1$ elsewhere, \m{|\delta|} its idele norm and $\chi$ the Hecke character attached to \m{\mathbf{a}_O} in
rank~\m{|\ix r|}.
\end{enumerate}
Let further \m{\ix N}, \m{\ix m}, \m{\ix a} be finite index sets, \m{z\colon\ix r\hookrightarrow\ix m} an embedding,
$\mathcal{V}$ a divided splitting in \m{\operatorname{Div}(\ix N,\ix m,\ix a;O,z)} and $s$ an \m{\ix a\times\ix N} matrix
over \m{\OEr/\mathfrak{M}_O^{2}}.
\begin{enumerate}[label=\textup{(\roman*)},leftmargin=2.3em,itemsep=2pt,resume]
\item \m{\stackOverGenScalar{\mathcal{V}}{s}} is the \m{(\ix a\sqcup\ix N)\times\ix N} matrix with $s$ stacked over
\m{\bar\varpi_{\mathcal{V}}\cdot1}.
\item \m{\congruenceGroup{\mathcal{V}}} is the group of the \m{A\in\GL_{\ix r}(\OEr/\mathfrak{M}_O^{2})} with
\m{A\cdot\bar\varpi_{\mathcal{V}}=\bar\varpi_{\mathcal{V}}}, that is \m{A\equiv1} modulo the chosen primes
(Input~M\,(a)).
\item A divided lattice~$L$ is \kw{glued for~$s$}, written \m{\IsGlued{\mathcal{V}}{s}{L}}, if (1) for every
\m{\ix n\times\ix r} matrix~$x$ over \m{\OEr/\mathfrak{M}_O^{2}} and every \m{c\in\CC},
\[\bigl[\reducedInverse{L}\,x=\bar\varpi_{\mathcal{V}}I_{e_O}\bigr]\,c=\sum_{A\in\congruenceGroup{\mathcal{V}}}
\bigl[(xA)_{\ix m}=I_z\ \text{and}\ \stackOverGenScalar{\mathcal{V}}{s}\,(xA)_{\ix N}=\stackOverGenScalar{\mathcal{V}}{s}\,I_{\phi_{\mathcal{V}}}\bigr]\,c,\]
where \m{(xA)_{\ix m}} and \m{(xA)_{\ix N}} are the rows \m{\varepsilon(\ix m)} and \m{\varepsilon(\ix N)} of~$xA$ and a
bracket is $1$ or~$0$; and (2) there is \m{K\in\GL_{\ix N}(\OEr/\mathfrak{M}_O^{2})} such that
\m{\bar\varpi_{\mathcal{V}}} times the columns \m{\varepsilon(\ix N)} of \m{\reducedInverse{L}}, multiplied by~$K$, vanishes
in every column outside the image of the divided embedding.
\item For a divided lattice~$L$, a pair \m{h=(h_1,h_2)} of points of \m{\UU(I_{\ix N})(\AAA_f)} and
\m{\UU(I_{\ix m})(\AAA_f)} and a coordinate vector~$y$ of weight~$\kappa^{+}_O$ with rows~$\ix N$, the \kw{product
form} is the linear map from coordinate vectors~$\eta$ of weight~$\kappa^{-}_O$ with rows~$\ix m$ to \m{\FormVx{O}}
\[\productForm{\mathcal{V}}{L}{h}{y}(\eta)=\pointForm{\operatorname{diag}_\varepsilon(h_1,h_2)\,\gamma_L}{I_{e_O}}{\varepsilon\circ\phi_{\mathcal{V}}}{y\cdot\eta},\]
where \m{\operatorname{diag}_\varepsilon} is the block diagonal embedding along~$\varepsilon$ and \m{y\cdot\eta}, the
\kw{Cartan product}, is the coordinate vector of weight~$\kappa_O$ obtained from \m{y\otimes\eta} by an integral
left inverse of the Cartan embedding.
\item The \kw{loss} of~$\mathcal{V}$ is
\m{\loss{\mathcal{V}}=\sum_\tau\bigl(|\ix a|\sum_{h<|\ix N|}\kappa^{+}_{O,\tau}(h)+|\ix m|\sum_{h<|\ix m|}\kappa^{-}_{O,\tau}(h)\bigr)}.
\end{enumerate}
\end{definition}
\leanline{\lean{DividedLattice}\sep \lean{reducedInverse}\sep \lean{weilScalar}\sep \lean{stackOverGenScalar}\sep
\lean{congruenceGroup}\sep \lean{IsGlued}\sep \lean{productForm}\sep \lean{loss}}

\begin{definition}[tables and comparison forms]\label{d:comp}
In Settings \ref{set:bg} and~\ref{set:data}, let $j$ be a lift index and $\vartheta$ a twist of~$D$.
\begin{enumerate}[label=\textup{(\roman*)},leftmargin=2.3em,itemsep=2pt]
\item \m{\smallHookNumber{j}} is the hook number of the weight~$\kappa^{-}$ of \m{O_D(1,j)} in the rank~\m{|\Six|}.
\item \m{\SmallFormValues{j}} is the space of coordinate vectors of weight~$\kappa^{-}$ with rows~$\Six$, with the
inner product of the Gram form; \m{\Tables{j}} is the orthogonal sum of copies of \m{\SmallFormValues{j}}, one for
each pair \m{a=(i,\mathfrak{c})} of a choice \m{i\in\Mem{j}} and a class \m{\mathfrak{c}\in\Scl}. Its elements, the
\kw{tables}, are families \m{\zeta=(\zeta_a)}.
\item \m{\splittingV{j}{\vartheta}} is the divided splitting, in
\m{\operatorname{Div}(\Bix{j},\Six,\Cix;O_D(\vartheta,j),z)} with $z$ the inclusion of~$\ix r$ in~$\Six$, given by the fixed
bijection \m{\Bix{j}\sqcup\Six\to\ix n}, the inclusion of~$\ix r$ in~$\Bix{j}$ as fixed embedding and a chosen
embedding of~$\Cix$ in~$\Bix{j}$ as divided embedding.
\item For a choice \m{i\in\Mem{j}}, \m{\glueLattice{j}{\vartheta}{i}} is a chosen divided lattice of
\m{O_D(\vartheta,j)} glued for~$i_1$ and \m{\splittingV{j}{\vartheta}}; its point is the one chosen for the trivial
twist, so it does not depend on~$\vartheta$.
\item \m{\lineScale{j}} is a chosen complex number \m{\ne0} with the properties of Input~L\,(c).
\item The \kw{comparison form} of a table~$\zeta$ is
\begin{align*}
\comparisonForm{j}{\vartheta}{\zeta}&=g(D)^{\loss{\splittingV{j}{\vartheta}}}\sum_{a=(i,\mathfrak{c})}\ \clsum{h}{\Gamma_i}
\productForm{\splittingV{j}{\vartheta}}{\glueLattice{j}{\vartheta}{i}}{(h,h_{\mathfrak{c}})}{y_i(h)}(\zeta_a),\\
y_i(h)&=\overline{\lineScale{j}\cdot\lmult{\vartheta}{j}{i}{h}},\qquad
\Gamma_i=\rowFixingLevel{O_D(\vartheta,j)}{\stackOverGenScalar{\splittingV{j}{\vartheta}}{i_1}}{\phi},
\end{align*}
with $\phi$ the fixed embedding. It is linear in~$\zeta$ and lies in \m{\FormVx{O_D(\vartheta,j)}}.
\end{enumerate}
\end{definition}
\leanline{\lean{smallHookNumber}\sep \lean{SmallFormValues}\sep \lean{Tables}\sep \lean{splittingV}\sep
\lean{glueLattice}\sep \lean{lineScale}\sep \lean{comparisonForm}\sep the three instances that make
\m{\SmallFormValues{j}} a finite-dimensional inner product space are unnamed}

\inputbox{\textbf{Inputs from the other notes of this unit.}

\vspace{4pt}\textbf{Input L} (note \emph{The height of a ratio of lift volumes}).\enspace In
Setting~\ref{set:data}, for a lift index~$j$:
(a) there are a real number~$C$ and numbers \m{c(\vartheta)\in\CC}, for the twists~$\vartheta$ of~$D$, with
\m{|c(\vartheta)|^2=C/\HV{j}{\vartheta}} such that for every twist~$\vartheta$, every choice \m{i\in\Mem{j}}, every
point~$g$ and all matrices $u$, $w$ over \m{\OEr/\bar\pi_g^{2}} with \m{uw=i_1}:
\m{\fnat{j}{i}(g)=c(\vartheta)\cdot\cks{O_D(\vartheta,j)}{\mathbf{b}''_D(\vartheta,j)}{w}{\phi}{\lmult{\vartheta}{j}{i}{\cdot}}(g)};
(b) a good form \m{F\in\Mgood{j}} with
\m{\int_{[\UU(W_j(D))]}\overline{F(g)\,\fnat{j}{i}(g)}\,\ttheta{j}{\zeta}(g)\,dg=0} for all choices~$i$ and all~$\zeta$ is
zero; for every good~$F$ and choice~$i$ the product \m{F\cdot\fnat{j}{i}} decreases rapidly, so that these integrals
converge;
(c) there is \m{c\ne0} in~$\CC$ such that for every choice~$i$, twist~$\vartheta$ and point~$h$ the vector
\m{y=\overline{c\cdot\lmult{\vartheta}{j}{i}{h}}} has gauge at most~$1$ at~$h$ for every prime~$p$ prime to the norms of
the primes in \m{Q(\vartheta)} and every \m{\iota_p}, and, if \m{Q(\vartheta)=\emptyset}, \m{\|\sigma(y)\|^2\le\exp(r\,B_0(D))}
for all \m{\sigma\in\AutC};
(d) \m{6r\le|\ix n|}-type bounds for the ranks: \m{|\Bix{j}|=r+9m+1}, \m{100m+9<k(D)}, \m{r+1=k(D)+m}.

\vspace{3pt}\textbf{Input M} (note \emph{Auxiliary dual pairs, determinant polynomials and the carrier}).\enspace
(a) For a divided splitting~$\mathcal{V}$ and a matrix~$S$ over \m{\OEr/\mathfrak{M}_O^{2}\mathfrak{Q}_O}:
\m{\bar\varpi_{\mathcal{V}}\operatorname{red}_0(S)=\bar\varpi_{\mathcal{V}}I_u} if and only if \m{S\equiv I_u} modulo every
chosen place. (b) An integral adele~$x$ reduces to~$0$ modulo \m{\mathfrak{M}_O^{c}} if and only if
\m{\varpi_O^{-c}x} is integral. (c) The descent \m{\dsc{O}} is the class sum of the conjugate kernels of the test datum
\m{\Upsilon_O} with the harmonic polynomial. (d) The bookkeeping lemmas on chosen lifts and block matrices.

\vspace{3pt}\textbf{Input V} (note \emph{The volume of the chain of theta lifts}).\enspace The adjoint lift commutes
with the actions of the group algebra of \m{\mathsf{G}(D)}.}

\inputbox{\textbf{Inputs from other units and from the background library.} Each is used where it is named in the
proofs; the notes on its unit, or the catalogue of the handout, state it in full.

\vspace{4pt}\textbf{Unit \emph{Forms}.}\enspace The conjugate kernel of an auxiliary pair is equivariant: at
\m{(g,\gamma xk)}, applied to \m{\chi_O(k)\rho(\gamma)v}, it equals its value at \m{(g,x)} applied to~$v$, for rational
$\gamma$ and $k$ in the level. The descent of a form is a multiple of the conjugate of a theta lift of its
conjugate.
\textbf{Unit \emph{SiegelWeil}.}\enspace (a) If \m{\psi_O\ne1} at each twist place, the reflection sums of all non-zero
rows at the chosen places kill~$\Xi$ and \m{\dsc{O}(\Xi)=0}, then the descents of the matrix coefficients of~$\Xi$ from
\m{\UU(I_{\ix n})} to \m{\UU(W_O)} vanish for data at the level. (b) If the descents of a function to~$\UU(W)$ vanish
and it is killed by the Hecke polynomial at a split place at the test value, then its descents to
\m{\UU(W\oplus\mathbb{H}^{p})} vanish, when \m{|\ix n|>2|\ix r|+2p+3}. (c) The representation of the finite group on a
level vector is translation by a lift times a character value.
\textbf{Unit \emph{Exchange}.}\enspace The pairs \m{O_D(\vartheta,j)} fit the parameter of the forms; the adjoint
lift is adjoint to the descent; the datum of a function modulo~$\mathfrak{m}$, integral away from~$q$, times a
power of an integer~$N$ with \m{qN\in\mathfrak{m}}, is a standard datum at primes not dividing~$q$, and it is an
admissible combination of translates of the vacuum of size \m{|\OEr/\mathfrak{m}|^{|\ix n||\ix n'|}} times a bound for the
conjugates of the function; \m{|\OEr/\mathfrak{m}|\le k^{2f}} if \m{k\in\mathfrak{m}}; the trivial lift of an $L$-value
datum and its normalised value.
\textbf{Unit \emph{Curve}.}\enspace Integrality of jets of standard data; conjugates of jets under
\m{\operatorname{Gal}(\CC/K)}, with the semilinear maps and their bounds; tight frames and the finite data they
produce; the limit to Katz's divisibility; Frobenius packings and large splits; a non-zero value of the integral
vacuum; the period factor.
\textbf{Units \emph{Data}, \emph{Columns}, \emph{Trace}, \emph{Seesaw}.}\enspace Levels inside the chosen level in
rank at most a fixed bound are free, and the orders of class stabilisers divide the chosen group order; the number
of classes of a definite unitary group at a level containing a principal level is at most an exponential in the
scales; comparison of monomials in the scales; \m{e_0(D)} is a central projection whose coefficients are algebraic
integers up to the order of the group, and it kills the reflection sums; the chain of lifts intertwines the finite
group with a right translation; \m{\dim\operatorname{Row}(D)=\dim\Mgood{j}\ne0}.
\textbf{Background library.}\enspace Fourier coefficients of theta lifts as periods: a continuous automorphic
function on \m{\UU(I_{\ix n})}, \m{\ix n=\ix n_1\sqcup\ix n_2}, fixed by a level, all of whose theta lifts to the
complementary group vanish, has vanishing period over \m{\UU(I_{\ix n_2})} (the catalogue gives P.~Walls and
S.~Rallis as references, place not confirmed, relation: a weaker form); the Hermite normal form over a valuation
ring; the Schrödinger-model formula at split places; the seesaw for \m{I_{\ix n}=I_{\ix n_1}\oplus I_{\ix n_2}}; the
translation rules of kernels; Schur matrices are multiplicative and contract the sup norm on integral matrices; the
Cartan product and the mixed Cartan embedding; duality for the Gram norm; the gauge of a delta form is at most the
gauge of its value; a finite Fourier expansion of a translated vacuum; the $p$-adic logarithm of the theta function
of a CM curve (Kronecker's first limit formula); references to be confirmed where the catalogue has none.}

\sect{Proofs}{s:proofs}
\subsection*{Point forms}
\begin{lemma}[a point with prescribed reductions]\label{l:reduce}
In Setting~\ref{set:bg}, let $O$ be an auxiliary pair for \m{(E,\ix r,\ix n)}, \m{Y=(Y_w)} a family of
\m{\ix n\times\ix r} matrices over the residue fields of the chosen places with linearly independent columns, and
\m{e\colon\ix r\hookrightarrow\ix n} an embedding. Then there is \m{k\in\KK{\ix n}}, with component~$1$ outside the
places of~$O$ and their conjugates, such that \m{(k\bmod w)\,I_{e_O}=Y_w} at every chosen place~$w$ and
\m{(k\bmod w)\transp I_e=I_{e_O}} at every twist place~$w$.
\end{lemma}
\begin{proof}
A matrix with independent columns over a field extends to an invertible matrix with these columns at the indices
\m{e_O(\ix r)}; this gives \m{A_w} with \m{A_wI_{e_O}=Y_w} at the chosen places and, applied to \m{I_{e_O}} with the
embedding~$e$, matrices~$B_w$ with \m{B_wI_e=I_{e_O}} at the twist places. The reduction from the elements of
\m{\KK{\ix n}} supported at the places of~$O$ to the product of the groups \m{\GL_{\ix n}(k_w)} is surjective; take
$k$ with reductions~$A_w$ and~$B_w\transp$.
\end{proof}
\leanline{\lean{existsreduce}}

\begin{lemma}[pairing with a point form]\label{l:inner}
In Setting~\ref{set:bg}, let $O$ be an auxiliary pair for \m{(E,\ix r,\ix n)}, \m{Y=(Y_w)} a family of matrices with
linearly independent columns and~$e$ an embedding as in Lemma~\ref{l:reduce}, and \m{h\in\UU(I_{\ix n})(\AAA_f)}. For every
form \m{\phi\in\FormVx{O}} and every coordinate vector~$v$ of weight~$\kappa_O$,
\m{\ip{\pointForm{h}{Y}{e}{v}}{\phi}=\ip{v}{\phi(h\,k_{Y,e})}_{\kappa_O}}.
\end{lemma}
\begin{proof}
The inner product is the class sum of the Gram forms of the values. The Gram form is invariant under the weight
representation of the rational points, which are unitary at infinity, and the character of the level has absolute
value~$1$. So the summand is constant on the class of \m{hk_{Y,e}}, where the delta form has the value~$v$, and it
vanishes on the other classes.
\end{proof}
\leanline{\lean{innerpointForm}}

\begin{lemma}[the descent of a point form]\label{l:thetapoint}
In Setting~\ref{set:bg}, let $O$ be an auxiliary pair for \m{(E,\ix r,\ix n)}, \m{Y=(Y_w)} a family of matrices with
linearly independent columns, $e$ an embedding and~$h$ a point as in Lemma~\ref{l:inner}, $v$ a coordinate vector of
weight~$\kappa_O$ and \m{g\in\UU(W_O)(\AAA)}. Then
\[\dsc{O}\bigl(\pointForm{h}{Y}{e}{v}\bigr)(g)=\overline{\tk{W_O}{I_{\ix n}}{\mathbf{a}_O}{\mathbf{b}_O}{\mathfrak{M}_O^{2}\mathfrak{Q}_O}{\Psi_{Y,e}}{P^{O}(\bar v)}(g,h)},\]
where \m{\Psi_{Y,e}(T)} is the product over the twist places~$w$ of \m{\psi_{O,w}(\det(I_e\transp T\bmod w))} if
\m{T\equiv Y_w} at every chosen place~$w$, and $0$ otherwise. The formal statement names a further character~$b$,
which does not occur in it.
\end{lemma}
\begin{proof}
The descent is the class sum of the conjugate kernels at \m{(g,\cdot)} applied to the values (Input~M\,(c)). By the
equivariance of the conjugate kernel (unit \emph{Forms}) the class sum for a delta form at \m{hk_{Y,e}} is the
kernel at \m{(g,hk_{Y,e})} applied to~$v$. The integral point \m{k_{Y,e}} is moved into the finite datum, which
becomes the translate of the datum of \m{\Upsilon_O} by the reductions of~\m{k_{Y,e}}; by the two properties of
Lemma~\ref{l:reduce} this translate is~\m{\Psi_{Y,e}}.
\end{proof}
\leanline{\lean{thetapointForm}}

\subsection*{Glued lattices}
\begin{lemma}[conjugating by a divided lattice]\label{l:conj}
In Setting~\ref{set:bg}, let $O$ be an auxiliary pair for \m{(E,\ix r,\ix n)}, \m{L\in\DividedLattice{O}} and
\m{j\in\KK{\ix n}} with \m{j\equiv1} modulo \m{\mathfrak{M}_O^{2}}. Then \m{\gamma_L^{-1}j\gamma_L} lies in \m{\KK{\ix n}} and
reduces to~$1$ at every chosen place of~$O$.
\end{lemma}
\begin{proof}
Write \m{\gamma=\gamma_L} and \m{\varpi=\varpi_O}. Since $\gamma$ is unitary, \m{\gamma=(\gamma^{-1})\herm}, so
\m{\bar\varpi\gamma} is integral together with \m{\varpi\gamma^{-1}}. The entries of \m{\varpi^{-2}(j-1)} are integral
by Input~M\,(b), and so are those of \m{\bar\varpi^{-1}\varpi^{-2}(j-1)}: at a chosen place \m{\bar\varpi} is a unit,
at the conjugate of a chosen place \m{\varpi=1} and \m{j-1} is the conjugate transpose of \m{j^{-1}-1}, which is
divisible by~$\varpi^{2}$ at the chosen place. Hence
\m{\varpi^{-1}(\gamma^{-1}j\gamma-1)=(\varpi\gamma^{-1})\cdot\bar\varpi^{-1}\varpi^{-2}(j-1)\cdot(\bar\varpi\gamma)} is
integral. So \m{\gamma^{-1}j\gamma} is integral, hence in~\m{\KK{\ix n}}, and it is~$1$ modulo each chosen prime.
\end{proof}
\leanline{\lean{reduceAtconj}}

\begin{lemma}[a divided lattice from a factorisation]\label{l:latexists}
In Setting~\ref{set:bg}, let $O$ be an auxiliary pair for \m{(E,\ix r,\ix n)} and let $C$, $G$ be \m{\ix n\times\ix n}
matrices over~$\Ohat$ with \m{CG=\varpi_O\cdot1}. Then there is \m{L\in\DividedLattice{O}} with
\m{\reducedInverse{L}=G} modulo \m{\mathfrak{M}_O^{2}}.
\end{lemma}
\begin{proof}
Let $c$ be the element of \m{\GL_{\ix n}(\AAA_f)} equal to~$C$ at the chosen places, with inverse
\m{\varpi_O^{-1}G} there, and to~$1$ elsewhere. Its conjugate transpose~$c\herm$ is supported at the conjugates of the
chosen places, a disjoint set, so $c$ and~$c\herm$ commute and \m{\gamma=(c\herm)^{-1}c} is unitary. At a chosen place
\m{\varpi_O\gamma^{-1}=G}; at its conjugate \m{\varpi_O=1} and \m{\gamma^{-1}=c\herm} is integral; elsewhere
\m{\gamma=1}.
\end{proof}
\leanline{\lean{existsreducedInverse}\sep the formal statement lists a lattice~$L$ among its arguments, which it
then binds itself}

\begin{lemma}[the levels of a divided splitting]\label{l:rowfix}
In Setting~\ref{set:bg}, let $O$ be an auxiliary pair for \m{(E,\ix r,\ix n)}, \m{\ix N}, \m{\ix m}, \m{\ix a} finite
index sets, \m{z\colon\ix r\hookrightarrow\ix m}, $\mathcal{V}$ a divided splitting in
\m{\operatorname{Div}(\ix N,\ix m,\ix a;O,z)} and $s$ an \m{\ix a\times\ix N} matrix over \m{\OEr/\mathfrak{M}_O^{2}}.
\begin{enumerate}[label=\textup{(\alph*)},leftmargin=2.3em,itemsep=2pt]
\item Every element of \m{\rowFixingLevel{O}{\stackOverGenScalar{\mathcal{V}}{s}}{\phi_{\mathcal{V}}}} reduces to~$1$ at every
chosen place of~$O$.
\item If $Y$ is an \m{\ix r\times\ix r} matrix over \m{\OEr/\mathfrak{M}_O^{2}} with \m{\bar\varpi_{\mathcal{V}}Y=0}, then
there is \m{A\in\congruenceGroup{\mathcal{V}}} with \m{A^{-1}=1-Y}.
\end{enumerate}
\end{lemma}
\begin{proof}
(a) The reduction~$A$ modulo \m{\mathfrak{M}_O^{2}} of such an element fixes the stacked matrix, in particular its
lower block: \m{\bar\varpi_{\mathcal{V}}A=\bar\varpi_{\mathcal{V}}\cdot1}. By Input~M\,(a), \m{A\equiv1} at every chosen
place. (b) The entries of~$Y$ lie in the image of~\m{\mathfrak{M}_O}, so \m{Y^2=0}, the matrix \m{1-Y} is invertible
and \m{\bar\varpi_{\mathcal{V}}(1-Y)=\bar\varpi_{\mathcal{V}}}; so \m{1-Y} and its inverse lie in
\m{\congruenceGroup{\mathcal{V}}}.
\end{proof}
\leanline{\lean{rowFixingLevelle}\sep \lean{existscongruenceGroupinv}}

\begin{remark}\label{r:liftC}
By Input~M\,(a), the chosen lift~$\tilde A$ of an element \m{A\in\congruenceGroup{\mathcal{V}}} reduces to~$1$ at every
chosen place. For the data, the level \m{\Gamma_i} of Definition~\ref{d:comp}\,(vi) is \kw{free}: the stabiliser in
the rational unitary group of each of its classes is trivial, by part~(a) of Lemma~\ref{l:rowfix}, the rank bounds
of Input~L\,(d) and the unit \emph{Data}. So a class sum at this level is a plain sum over the classes.
\end{remark}
\leanline{\lean{reduceAtliftMod}\sep \lean{classAutrowFixingLevel}}

\begin{proposition}[glued lattices exist]\label{p:glue}
In Setting~\ref{set:bg}, let $O$ be an auxiliary pair for \m{(E,\ix r,\ix n)}, \m{\ix N}, \m{\ix m}, \m{\ix a} finite
index sets, \m{z\colon\ix r\hookrightarrow\ix m}, $\mathcal{V}$ a divided splitting in
\m{\operatorname{Div}(\ix N,\ix m,\ix a;O,z)} and $s$ an \m{\ix a\times\ix N} matrix over \m{\OEr/\mathfrak{M}_O^{2}}. Then
there is a divided lattice \m{L\in\DividedLattice{O}} with \m{\IsGlued{\mathcal{V}}{s}{L}}.
\end{proposition}
\begin{proof}
Lift $s$ to a matrix~$t$ over~$\Ohat$. By the Hermite normal form over~$\Ohat$ (library), the matrix with $t$
stacked over \m{\varpi_O\cdot1} is \m{U_{\ix N}D}, where $U$ is invertible over~$\Ohat$, $U_{\ix N}$ are its last
columns and $D$ is a square matrix; the lower block gives \m{U_{22}D=\varpi_O\cdot1}. Along~$\varepsilon$, let
\[C=\begin{pmatrix}U_{22}&I_{\phi}I_z\transp\\0&\Delta\end{pmatrix},\qquad
G=\begin{pmatrix}D&-D\,I_{\phi}I_z\transp\\0&\Delta'\end{pmatrix},\]
with $\phi$ the fixed embedding, $\Delta$ the diagonal matrix with entries~$1$ on \m{z(\ix r)} and \m{\varpi_O}
elsewhere, and $\Delta'$ the one with \m{\varpi_O} on \m{z(\ix r)} and $1$ elsewhere. Then \m{CG=\varpi_O\cdot1}, and
Lemma~\ref{l:latexists} gives~$L$ with \m{\reducedInverse{L}=G} modulo \m{\mathfrak{M}_O^{2}}.
For condition~(1): the lower block of \m{\reducedInverse{L}\,x=\bar\varpi_{\mathcal{V}}I_{e_O}} says that the rows
\m{\varepsilon(\ix m)} of~$x$ outside \m{z(\ix r)} vanish and that the rows \m{z(\ix r)} form a matrix \m{1-Y} with
\m{\bar\varpi_{\mathcal{V}}Y=0}; by Lemma~\ref{l:rowfix}\,(b) there is exactly one \m{A\in\congruenceGroup{\mathcal{V}}}
with \m{(xA)_{\ix m}=I_z}. For this~$A$ the upper block says \m{\bar D\,(xA)_{\ix N}=\bar D\,I_\phi}, with $\bar D$ the
reduction of~$D$, and this is equivalent to the same equation with the stacked matrix in the place of~$\bar D$,
because the stacked matrix is \m{\bar U_{\ix N}\bar D} and $U$ is invertible. For condition~(2): the Hermite
normal form over the residue ring gives an invertible~$K$ such that $sK$ vanishes outside the columns of the
divided embedding; then \m{\bar\varpi_{\mathcal{V}}\bar DK} vanishes there too, since $\bar D$ is a combination of the
rows of~$s$ and of \m{\bar\varpi_{\mathcal{V}}\cdot1}, and \m{\bar\varpi_{\mathcal{V}}^{2}=0}.
\end{proof}
\leanline{\lean{existsisGlued}\sep the formal statement lists a lattice~$L$ among its arguments, which it binds
itself}

\subsection*{The descent of a product form}
\begin{lemma}[the descent of a point form at a divided lattice]\label{l:thetaL}
In Setting~\ref{set:bg}, let $O$ be an auxiliary pair for \m{(E,\ix r,\ix n)}, \m{\ix N}, \m{\ix m} finite index sets,
\m{\ix a} an index set, \m{z\colon\ix r\hookrightarrow\ix m}, $\mathcal{V}$ a divided splitting in
\m{\operatorname{Div}(\ix N,\ix m,\ix a;O,z)}, \m{L\in\DividedLattice{O}}, \m{e\colon\ix r\hookrightarrow\ix n} an embedding,
$v$ a coordinate vector of weight~$\kappa_O$, \m{g\in\UU(W_O)(\AAA)} and \m{h\in\UU(I_{\ix n})(\AAA_f)}. Then
\[\dsc{O}\bigl(\pointForm{h\gamma_L}{I_{e_O}}{e}{v}\bigr)(g)=\overline{\weilScalar{L}\cdot\tk{W_O}{I_{\ix n}}{\mathbf{a}_O}{\mathbf{b}_O}{\mathfrak{M}_O^{2}\mathfrak{Q}_O}{\Psi_{L,e}}{P^{O}(\bar v)}(g,h)},\]
where \m{\Psi_{L,e}(S)} is the product over the twist places~$w$ of the numbers
\m{\psi_{O,w}(\det(I_e\transp S\bmod w))} if
\[\reducedInverse{L}\operatorname{red}_0(S)=\bar\varpi_{\mathcal{V}}I_{e_O},\]
and $0$ otherwise. The formal statement names
a character~$b$, which does not occur in it.
\end{lemma}
\begin{proof}
By Lemma~\ref{l:thetapoint} at the point \m{h\gamma_L}. The point~$\gamma_L$ is supported at the chosen places and
their conjugates, where the unitary group is a general linear group, and the Schrödinger-model formula (library)
moves it into the datum: the function \m{x\mapsto\Psi(x)} becomes \m{x\mapsto\Psi(\gamma_L^{-1}x)}, with the factor
\m{\chi(\delta)|\delta|^{-|\ix r|/2}=\weilScalar{L}}. The condition ``\m{\gamma_L^{-1}x} is integral and congruent to
\m{I_{e_O}} at the chosen places'' is, by Input~M\,(a) and~(b), the condition on \m{\reducedInverse{L}}.
\end{proof}
\leanline{\lean{thetapointFormL}}

\begin{proposition}[the descent of a product form]\label{p:thetaprod}
In Setting~\ref{set:bg}, let $O$ be an auxiliary pair for \m{(E,\ix r,\ix n)}, \m{\ix N}, \m{\ix m}, \m{\ix a} finite
index sets, \m{z\colon\ix r\hookrightarrow\ix m}, $\mathcal{V}$ a divided splitting in
\m{\operatorname{Div}(\ix N,\ix m,\ix a;O,z)}, $s$ an \m{\ix a\times\ix N} matrix over \m{\OEr/\mathfrak{M}_O^{2}},
\m{L\in\DividedLattice{O}} with \m{\IsGlued{\mathcal{V}}{s}{L}}, and $b$, $b'$ characters with \m{b\,b'=\mathbf{b}_O}.
Write \m{w_s=\stackOverGenScalar{\mathcal{V}}{s}}, \m{\phi=\phi_{\mathcal{V}}}, and
\m{\Theta^{-}_{\eta}(g,h')=\tk{W_O}{I_{\ix m}}{\mathbf{a}_O}{b'}{\mathfrak{M}_O^{2}}{\mathbf{1}_{I_z}}{\operatorname{pol}^{-}_{O,\kappa^{-}_O}(\bar\eta)}(g,h')}.
\begin{enumerate}[label=\textup{(\alph*)},leftmargin=2.3em,itemsep=2pt]
\item For points $h$ of \m{\UU(I_{\ix N})(\AAA_f)} and $h'$ of \m{\UU(I_{\ix m})(\AAA_f)}, coordinate vectors $y$ of
weight~$\kappa^{+}_O$ and $\eta$ of weight~$\kappa^{-}_O$ and \m{g\in\UU(W_O)(\AAA)}:
\begin{align*}
\dsc{O}\bigl(\productForm{\mathcal{V}}{L}{(h,h')}{y}(\eta)\bigr)(g)&=\overline{\weilScalar{L}\sum_{A\in\congruenceGroup{\mathcal{V}}}
\Theta^{+}_{y}(g\tilde A,h)\cdot\Theta^{-}_{\eta}(g\tilde A,h')},\\
\Theta^{+}_{y}(g,h)&=\tk{W_O}{I_{\ix N}}{\mathbf{a}_O}{b}{\mathfrak{M}_O^{2}\mathfrak{Q}_O}{\operatorname{cd}_{\chi,\operatorname{red}_0,w_s,\phi}}{\operatorname{pol}^{+}_{O,\kappa^{+}_O}(\bar y)}(g,h).
\end{align*}
\item For a point~$h'$, \m{c\in\CC}, a function~$y$ from \m{\UU(I_{\ix N})(\AAA_f)} to coordinate vectors of
weight~$\kappa^{+}_O$, a vector~$\eta$ and~$g$ as before:
\begin{align*}
&\dsc{O}\Bigl(\clsum{h}{\rowFixingLevel{O}{w_s}{\phi}}\productForm{\mathcal{V}}{L}{(h,h')}{\overline{c\,y(h)}}(\eta)\Bigr)(g)\\
&\qquad=\overline{\weilScalar{L}\,c}\cdot\sum_{A\in\congruenceGroup{\mathcal{V}}}\overline{\cks{O}{b}{w_s}{\phi}{y}(g\tilde A)\cdot\Theta^{-}_{\eta}(g\tilde A,h')}.
\end{align*}
\end{enumerate}
\end{proposition}
\begin{proof}
(a) Apply Lemma~\ref{l:thetaL} to the point \m{\operatorname{diag}_\varepsilon(h,h')} and the vector \m{y\cdot\eta}. By
condition~(1) of a glued lattice the datum \m{\Psi_{L,e}} is the sum over \m{A\in\congruenceGroup{\mathcal{V}}} of the
translates by~$A$ of the product of the indicator of \m{I_z} on the rows~$\ix m$ and of the column datum on the
rows~$\ix N$; the factors \m{\psi_{O,w}} at the twist places are the character~$\chi$ of the column datum, because
\m{\tilde A} is~$1$ there. The translation by~$A$ is moved to the first variable (library), which gives
\m{g\tilde A}. The harmonic polynomial of a Cartan product is the product of the polynomial of~$y$ in the positive
letters and of~$\eta$ in the negative letters, the columns~$\ix N$ and~$\ix m$ carrying no columns of the other
weight by the conditions of a divided splitting (library: the mixed Cartan embedding). The seesaw for
\m{I_{\ix n}=I_{\ix N}\oplus I_{\ix m}}, with \m{\mathbf{b}_O=b\,b'}, splits the kernel into the product, and the second
factor only depends on its datum modulo \m{\mathfrak{M}_O^{2}}.
(b) The descent is linear and commutes with the finite class sum; apply~(a) to each term and use that kernels are
linear in the polynomial.
\end{proof}
\leanline{\lean{thetaproductForm}\sep \lean{thetastdClassSumproductForm}}

\subsection*{Forms that are not lifts have vanishing theta periods}
\begin{setting}[a form that is not a lift]\label{set:nonlift}
In Setting~\ref{set:bg}: \m{\ix r} is not empty, \m{\ix N} is a finite index set, $O$ is an auxiliary pair for
\m{(E,\ix r,\ix n)} that fits a parameter~$l_0$ of forms on~$W_O$ in the sense of the unit \emph{Exchange}, with
\m{\psi_{O,w}\ne1} at every twist place~$w$ and \m{2|\ix r|+2|\ix N|+3<|\ix n|}. $\Xi$ is a form in \m{\FormVx{O}}
such that \m{\dsc{O}(\Xi)=0}; such that for every chosen place~$w_0$ and every family~$v$ of rows over the residue
fields of the chosen places with \m{v_{w_0}\ne0} the tensor product of the reflection sums of \m{\psi_O} and~$v$ acts
on~$\Xi$ by zero; and such that \m{H_{O,w}(z_O(w))\,\Xi=0} for a finite place~$w$ with \m{w\ne\bar w} and $w$,
$\bar w$ not among the places of~$O$.
\end{setting}

\begin{lemma}[descents to a sum with hyperbolic planes]\label{l:hyp}
In Setting~\ref{set:nonlift}, for every coordinate vector~$u$ of weight~$\kappa_O$, the complex conjugate of the
matrix coefficient \m{h\mapsto\ip{u}{\Xi(h)}_{\kappa_O}} has vanishing theta lifts from \m{\UU(I_{\ix n})} to the unitary
group of \m{W_O\oplus\mathbb{H}^{|\ix N|}}, the sum of~$W_O$ with \m{|\ix N|} hyperbolic planes, for the character
\m{\mathbf{a}_O\upsilon^{|\ix N|}} and all data that are the vacuum outside the places of~$O$ and their conjugates.
\end{lemma}
\begin{proof}
By the unit \emph{SiegelWeil}\,(a) the descents of the matrix coefficients to \m{\UU(W_O)} vanish: its hypotheses
are those of Setting~\ref{set:nonlift}, the vanishing of \m{\dsc{O}(\Xi)} being translated by the unit \emph{Forms}
into the vanishing of a theta lift of the conjugate coefficient against the highest weight vector, and extended to
all~$u$ because the highest weight vector is cyclic. By the unit \emph{SiegelWeil}\,(b), with \m{p=|\ix N|}, the
descents to \m{\UU(W_O\oplus\mathbb{H}^{|\ix N|})} vanish: the coefficient is spherical at~$w$ and killed by the Hecke
polynomial at the test value, since $\Xi$ is.
\end{proof}
\leanline{\lean{descentsVanishhypSum}}

\begin{proposition}[vanishing periods]\label{p:period}
In Setting~\ref{set:nonlift}, let \m{\ix l} be a finite index set, \m{\varepsilon\colon\ix N\sqcup\ix l\to\ix n} a
bijection, $\kappa$ a weight, \m{Y=(Y_w)} matrices with independent columns as in Lemma~\ref{l:reduce}, and
\m{e'\colon\ix r\hookrightarrow\ix n} an embedding of the form \m{\varepsilon\circ u} with \m{u\colon\ix r\to\ix N}. Let
$C$ be a conjugate-linear map from the coordinate vectors of weight~$\kappa$ with rows~$\ix l$ to those of
weight~$\kappa_O$ that is equivariant for the definite group on~$\ix l$ embedded as the block
\m{\operatorname{diag}_\varepsilon(1,y)}. Let $b$ be a character, $\mathfrak{m}$ a modulus whose primes are places
of~$O$, $\Psi$ a function on matrices modulo~$\mathfrak{m}$, $g$ and~$h_1$ points, \m{x\in\UU(I_{\ix n})(\AAA_f)} a
point with component~$1$ outside the chosen places and their conjugates, $K$ a level of \m{\UU(I_{\ix l})}, and let
\[f(h)=\Bigl\langle\pointForm{\operatorname{diag}_\varepsilon(h_1,h)\,x}{Y}{e'}{C\bigl(\kernelVector{O}{b}{\mathfrak{m}}{\Psi}{\kappa}{g}{h}{\cdot}\bigr)},\ \Xi\Bigr\rangle .\]
Suppose that every \m{k\in K} satisfies \m{\Psi((k\bmod\mathfrak{m})^{-1}T)=\Psi(T)} for all~$T$, and that the
element \m{x^{-1}\operatorname{diag}_\varepsilon(1,k)\,x} lies in \m{\KK{\ix n}} and its reduction fixes every~$Y_w$. Then
\m{\clsum{h}{K}f(h)=0}.
\end{proposition}
\begin{proof}
By Lemma~\ref{l:inner}, \m{f(h)} is the Gram form of \m{C(\text{kernel vector at }(g,h))} with the value of~$\Xi$ at
\m{\operatorname{diag}_\varepsilon(h_1,h)\,x\,k_{Y,e'}}. Consider on \m{\UU(I_{\ix n})(\AAA)} the function
\[F(h)=\sum_{i}\theta\bigl(\omega(g,h)\varphi_i\bigr)\cdot\overline{\ip{C(\delta_i)}{\Xi(h\,x\,k_{Y,e'})}_{\kappa_O}},\]
where $i$ runs over the standard fillings of weight~$\kappa$, $\delta_i$ is the $i$-th basis vector, and
$\varphi_i$ is the product, for the decomposition \m{I_{\ix n}=I_{\ix N}\oplus I_{\ix l}}, of a fixed translate~$v$ of the
vacuum on the first summand and of the datum of~$\Psi$ with the $i$-th basis polynomial in the negative letters on
the second. It is left invariant under the rational points and right invariant under a level. On the block
\m{\operatorname{diag}_\varepsilon(h_1,(y,h))}, by the seesaw in the second space, by the equivariance of~$C$ and by the
rule for the real Weil representation on Fock vectors, \m{F=\theta(\omega(g,h_1)v)\cdot\overline{f(h)}}. The theta
lifts of~$F$ to the complementary group \m{\UU(W_O\oplus\mathbb{H}^{|\ix N|})} are, term by term, the lifts of
Lemma~\ref{l:hyp} translated by \m{x\,k_{Y,e'}}, the data \m{\varphi_i} being the vacuum off the places of~$O$; so they
vanish. By the theorem on periods (library), the period of~$F$ over \m{\UU(I_{\ix l})} vanishes:
\m{\theta(\omega(g,h_1)v)\cdot\overline{\int f}=0}. The translate~$v$ can be chosen with
\m{\theta(\omega(g,h_1)v)\ne0}, because the theta functional does not vanish on all translates of the vacuum. Finally
the integral of~$f$ over the definite group is its class sum at level~$K$, since $f$ is right invariant under~$K$:
a translation by \m{k\in K} is absorbed by~$\Psi$ in the kernel vector, by the hypothesis on~$\Psi$, and by the
level of~$\Xi$ in its value, by the hypothesis on \m{x^{-1}\operatorname{diag}_\varepsilon(1,k)x} and the unit
\emph{SiegelWeil}\,(c).
\end{proof}
\leanline{\lean{stdClassSumperiod}}

\begin{corollary}[product forms are orthogonal to forms that are not lifts]\label{c:orth}
In Setting~\ref{set:nonlift}, let \m{\ix l} be a finite index set, \m{\ix a} an index set,
\m{z\colon\ix r\hookrightarrow\ix l}, $\mathcal{V}$ a divided splitting in \m{\operatorname{Div}(\ix N,\ix l,\ix a;O,z)}, and
suppose that \m{\mathfrak{M}_O^{2}} is generated by an element \m{\varpi\ne0} of~$\OEr$. Let \m{L\in\DividedLattice{O}},
$x$ a point of \m{\UU(I_{\ix N})(\AAA_f)}, $y$ a coordinate vector of weight~$\kappa^{+}_O$, $b$ a character, $\Psi$ a
function on matrices modulo \m{\mathfrak{M}_O^{2}} and $g$ a point. Then, with \m{K(\varpi)} the principal level
of \m{\UU(I_{\ix l})} modulo~$\varpi$,
\[\clsum{h}{K(\varpi)}\Bigl\langle\productForm{\mathcal{V}}{L}{(x,h)}{y}\Bigl(\overline{\kernelVector{O}{b}{\mathfrak{M}_O^{2}}{\Psi}{\kappa^{-}_O}{g}{h}{\cdot}}\Bigr),\ \Xi\Bigr\rangle=0 .\]
\end{corollary}
\begin{proof}
This is Proposition~\ref{p:period} for \m{\varepsilon} of~$\mathcal{V}$, \m{Y_w=I_{e_O}},
\m{e'=\varepsilon\circ\phi_{\mathcal{V}}}, the point \m{\gamma_L}, the map \m{C(\zeta)=y\cdot\bar\zeta} and the level
\m{K(\varpi)}. The map~$C$ is equivariant because the weight representation of a block diagonal element on a Cartan
product is the Cartan product of the representations (library). An element~$k$ of \m{K(\varpi)} is~$1$ modulo
\m{\mathfrak{M}_O^{2}}, so it fixes~$\Psi$, and by Lemma~\ref{l:conj} the conjugate of
\m{\operatorname{diag}_\varepsilon(1,k)} by~$\gamma_L$ is integral and reduces to~$1$ at the chosen places.
\end{proof}
\leanline{\lean{stdClassSuminnerproductForm}}

\subsection*{Integrality of product forms}
\begin{lemma}[the loss of the Cartan product]\label{l:cartan}
In Setting~\ref{set:bg}, let $O$ be an auxiliary pair for \m{(E,\ix r,\ix n)}, \m{\ix N}, \m{\ix m}, \m{\ix l} finite
index sets, \m{\ix a} an index set, \m{z\colon\ix r\hookrightarrow\ix m}, $\mathcal{V}$ a divided splitting in
\m{\operatorname{Div}(\ix N,\ix m,\ix a;O,z)}, $p$ a prime with an embedding \m{\iota_p\colon\CC\to\CC_p}, and \m{M\ge1} an
integer in the ideal \m{\mathfrak{M}_O}.
\begin{enumerate}[label=\textup{(\alph*)},leftmargin=2.3em,itemsep=2pt]
\item Let \m{L\in\DividedLattice{O}}, $\kappa$ a weight without columns of height~$h$ for
\m{|\ix l|\le h<|\ix n|}, with \m{|\ix l|\le|\ix n|}, let \m{f\colon\ix l\to\ix n} be a map, $S$ a set of columns,
$g$ and~$h$ finite adelic points and $v$ a coordinate vector of weight~$\kappa$. Suppose that some invertible~$K$
over \m{\OEr/\mathfrak{M}_O^{2}} makes \m{\bar\varpi_{\mathcal{V}}} times the columns~$f$ of \m{\reducedInverse{L}}, multiplied
by~$K$, vanish outside the columns~$S$, and that the columns~$f$ of~$g^{-1}$ are \m{I_fh^{-1}}. Then the sup norm
of the Schur matrix of the columns~$f$ of \m{(g\gamma_L)^{-1}} applied to \m{\iota_p(v)} is at most
\m{|M|_p^{-\nu}\cdot\gau{\iota_p}{\kappa}{0}{h}{v}}, with \m{\nu=\sum_\tau|S|\sum_{i<|\ix l|}\kappa_\tau(i)}.
\item Let $s$ be an \m{\ix a\times\ix N} matrix over \m{\OEr/\mathfrak{M}_O^{2}}, with $\ix a$ finite,
\m{L\in\DividedLattice{O}} with \m{\IsGlued{\mathcal{V}}{s}{L}}, \m{h=(h_1,h_2)} a pair of points and $y$, $\eta$
coordinate vectors of weights $\kappa^{+}_O$ and~$\kappa^{-}_O$. Then
\[\gau{\iota_p}{\kappa_O}{\epsilon_O}{\operatorname{diag}_\varepsilon(h_1,h_2)\gamma_L}{y\cdot\eta}\le|M|_p^{-\loss{\mathcal{V}}}\cdot
\gau{\iota_p}{\kappa^{+}_O}{0}{h_1}{y}\cdot\gau{\iota_p\circ\mathrm{conj}}{\kappa^{-}_O}{0}{h_2}{\bar\eta}.\]
\end{enumerate}
\end{lemma}
\begin{proof}
(a) Lift~$K$ to an integral point~$k$; replacing $h$ by~$hk^{-1}$ does not change the gauge. The matrix
\m{\varpi_O\gamma_L^{-1}} is integral, so the columns~$f$ of \m{(g\gamma_L)^{-1}}, multiplied by~$k$ and by the diagonal
matrix with entries~$M$ on~$S$ and $1$ elsewhere, form an integral matrix: on~$S$ because \m{M\in\mathfrak{M}_O}, and
outside~$S$ by the hypothesis and Input~M\,(b). Schur matrices are multiplicative and contract the sup norm on
integral matrices (library), and the Schur matrix of the inverse diagonal matrix has norm at most
\m{|M|_p^{-\nu}}. (b) The weight representation of a block diagonal element on a Cartan product is the Cartan product
of the representations, and an integral left inverse of the Cartan embedding has norm at most~$1$. Apply (a) to the
two blocks: to \m{\kappa^{+}_O} on~$\ix N$ with $S$ the image of the divided embedding, by condition~(2) of a glued
lattice, and to \m{\kappa^{-}_O} on~$\ix m$ with $S$ everything, for the conjugate embedding, since the block
on~$\ix m$ enters through the inverse transpose. The two exponents add to \m{\loss{\mathcal{V}}}.
\end{proof}
\leanline{\lean{normschurinvcols}\sep \lean{weylGaugecartanProduct}}

\subsection*{Crude denominators and conjugate bounds for theta jets}
The next results concern pairs of \emph{definite} standard spaces over the fields of the parameters. The notions are
those of the units \emph{Data} and \emph{Curve}: for parameters \m{P=(N_0,d)} with \m{K=\mathbb{Q}(\sqrt{-d})} and an
integer \m{s\ge0}, $E_s$ is the CM field of the parameters with exponent~$s$ (for data, \m{E(D)} is such a field),
$\Phi_s$ its set of infinite places, \m{N_1(s)} its conductor, \m{\Omega_s>0} its period constant,
\m{\operatorname{per}_s} its jet scale, \m{\kappa_0(s)} its vacuum scale and $c_P$ a constant of~$P$. For a Schwartz
datum~$f$ of a pair of definite standard spaces and \m{c\in\CC}, \m{T_c(f)} is the \kw{scaled Taylor series} of~$f$: a
formal power series in variables \m{X_{\tau,i,i'}} whose coefficients are, up to the scale~$c$, the Taylor
coefficients at the origin of the theta function of~$f$ in the coordinates of the CM curve. The \kw{jet} of a power
series~$F$ along a polynomial~$Q$ is \m{\operatorname{jet}_Q(F)=Q(\partial)F(0)}, and \m{\operatorname{jet}^{\kappa}(F)} is the
matrix whose product with a coordinate vector~$\xi$ of weight~$\kappa$ has as coordinates the jets of~$F$ along
the coefficient polynomials of \m{z\mapsto S_\kappa(z)\xi}. An \kw{admissible combination of size~$B$} of translates of
the integral vacuum is a finite combination of Heisenberg translates of the integral vacuum by finite adelic vectors
whose coefficients have, under every element of \m{\operatorname{Gal}(\CC/K)}, absolute values of sum at most~$B$. The
\kw{boxes} of a weight~$\kappa$ below~$N$ are \m{\sum_\tau\sum_{h<N}h\,\kappa_\tau(h)}, and
\m{\lambda_\kappa(\tau,i)=\sum_{i<h<N}\kappa_\tau(h)} is its row length.

\begin{lemma}[small facts]\label{l:small}
\textup{(a)} For parameters~$P$ and every \m{\sigma\in\AutC} there is \m{\sigma'\in\operatorname{Gal}(\CC/K)} with
\m{|\sigma(z)|=|\sigma'(z)|} for all \m{z\in\CC}.
\textup{(b)} Let \m{\ix L} and~\m{\ix N} be finite index sets, \m{\ix L} not empty, $e$ a family of integers and $l$
a family of integer vectors, indexed by~$\ix N$ and by \m{\ix N\times\ix L}, and $\kappa$ a weight such that $l$ is
twice the Harish-Chandra parameter of \m{(\kappa,e)}. If \m{\beta\ge0} is an integer with \m{|l(w,i)|\le\beta} for all
$w$, $i$, then the boxes of~$\kappa$ below~\m{|\ix L|} are at most \m{|\ix N|\,|\ix L|\,\beta}, and \m{|e(w)|\le\beta}
for all~$w$.
\textup{(c)} Let $\Phi$ and~$\ix n$ be finite, $s$ a sign vector on~$\ix n$, $\kappa$ a weight, $e$ a family of
integers indexed by~$\Phi$, $c$ a real number and $l$ a family of integer vectors on~$\ix n$ indexed by~$\Phi$. If
every~$l_w$ has gap~$c$ for the signs~$s$, in the sense of the unit on discrete series, and
\m{e(w)\le l(w,i)/2+|\ix n|/2} for all $w$, $i$, with integer division, then there is a weight~$f$ such that $l$ is
twice the Harish-Chandra parameter of \m{(f,\,w\mapsto e(w)+f_w(|\ix n|))} and \m{f_w(k)=\kappa_w(k)} for all $w$ and
all \m{k>|\ix n|}.
\textup{(d)} In Setting~\ref{set:bg}, let \m{\ix L}, \m{\ix N} be finite index sets, $a$ and~$b$ characters with
weights $e$, $e'$ and tame parts $\psi$, $\psi'$, $\mathfrak{m}$ a modulus, $\Psi$ a function on matrices
modulo~$\mathfrak{m}$, $Q$ a polynomial, $\gamma$ and~$\gamma'$ rational points of \m{\UU(I_{\ix L})} and
\m{\UU(I_{\ix N})}, and $g$, $h$ finite adelic points. Then
\[\tk{I_{\ix L}}{I_{\ix N}}{a}{b}{\mathfrak{m}}{\Psi}{Q}(\gamma g,\gamma'h)=\prod_\tau\tau(\det\gamma')^{-e_\tau}\tau(\det\gamma)^{e'_\tau}\cdot
\tk{I_{\ix L}}{I_{\ix N}}{a}{b}{\mathfrak{m}}{\Psi}{Q'}(g,h),\]
where $\tau$ also denotes the embedding of the CM type at~$\tau$ and $Q'$ is $Q$ after the substitution
\m{x_\tau\mapsto\tau(\gamma'^{-1})\transp x_\tau\,\tau(\gamma)\transp}.
\end{lemma}
\begin{proof}
(a) The restriction of~$\sigma$ to the imaginary quadratic field~$K$ is the identity or complex conjugation; take
\m{\sigma'=\sigma} or \m{\sigma'=\sigma\circ\mathrm{conj}}. (b) The boxes are the sum of the row lengths over the
\m{|\ix N|\,|\ix L|} pairs \m{(w,i)}, and each row length and each \m{|e(w)|} is recovered from two values of~$l_w$
(library), hence at most~$\beta$. (c) By induction on the values of~$l_w$ in increasing order, defining $f_w$ at
one more height in each step so that the next value is matched; the gap makes the differences positive. (d) The
theta functional is invariant under rational points, and at the infinite places a rational point acts on the Fock
vector of~$Q$ through the real Weil representation, which substitutes the matrices into~$Q$ and multiplies by the
determinant characters (Input~6 of the sister note); the weights of $a$ and~$b$ give the exponents.
\end{proof}
\leanline{\lean{existsalgEquiv}\sep \lean{boxesle}\sep \lean{existsisHCParamofle}\sep \lean{modKernelembFin}\sep with
\m{\kappa=0} and $e$ the minimum of the values, (c) gives: every family with a gap is twice a Harish-Chandra
parameter, \lean{existsisHCParamofhasGap}\sep the two bookkeeping lemmas \lean{sumdiagramWeight} (the sum of the
row lengths is the number of boxes) and \lean{summemintegralAbove} (a block sum of points integral above~$p$ is
integral above~$p$) need no argument}

\begin{proposition}[conjugates of jets]\label{p:jet}
Let $P$ be parameters, \m{s\ge0} an integer, \m{\ix n} and~\m{\ix n'} finite index sets, $a$ and~$b$ characters of
\m{E_s} with weights $e_a$, $e_b$ and tame parts $\psi$, $\psi'$, $B$ a real number and $f$ a datum of the pair
\m{(I_{\ix n},I_{\ix n'})} that is an admissible combination of size~$B$ of translates of the integral vacuum. Let
\m{\sigma\in\AutC}, \m{c\in\CC}, and let $g$, $h$ be finite adelic points of the two unitary groups. Put
\m{u=\log(|\sigma(c)|^2/\Omega_s)},
\[M(L,w)=B\cdot\exp\bigl(c_P\,(N_1(s)\,|\ix n'|\,|\ix n|+1)^2+|\Phi_s|\,|\ix n'|\,|\ix n|\,\log(2L+2)+w\bigr),\]
and \m{F=T_{c\operatorname{per}_s}\bigl(\omega_{a,b}(g,h)f\bigr)}. Then:
\begin{enumerate}[label=\textup{(\roman*)},leftmargin=2.3em,itemsep=2pt]
\item for every polynomial~$Q$, every integer \m{L\ge0}, every real~$w$ and every real \m{A\ge0}: if
\m{|\sigma(Q(z))|\le A\,e^{w+LR}} for all points~$z$ and all real~$R$ with
\m{\sum_i|\sigma(z_i)|^2\le e^{u}(2L+2)R}, then \m{|\sigma(\operatorname{jet}_Q(F))|\le M(L,w)\cdot A};
\item for every integer \m{L\ge1}, every weight \m{(\kappa,e)} that is a lift weight for \m{(\ix n,\ix n',e_b)} and
every coordinate vector~$\xi$:
\begin{align*}
\|\sigma(\operatorname{jet}^{\kappa}(F)\xi)\|_\kappa&\le M(L,w_\kappa)\cdot\|\sigma(\xi)\|_\kappa,\\
w_\kappa&=\tfrac12\sum_\tau\sum_{i<|\ix n|}\Bigl(\frac{\lambda_\kappa(\tau,i)}{L}+\log(\lambda_\kappa(\tau,i)!)+\lambda_\kappa(\tau,i)\,u\Bigr).
\end{align*}
\end{enumerate}
\end{proposition}
\begin{proof}
(i) By Lemma~\ref{l:small}\,(a) one may take \m{\sigma\in\operatorname{Gal}(\CC/K)}. The datum is a combination of
translates of the vacuum with coefficients of total size~$B$ under~$\sigma$, so it suffices to bound each term.
For a translate of the vacuum, the conjugate of a jet is the jet of the image under a semilinear map, of norm at
most \m{\exp(c_P(N_1(s)|\ix n'||\ix n|+1)^2)} (unit \emph{Curve}), and a jet of a scaled Taylor series is the theta
functional of the datum integrated against the Fock vector of the scaled polynomial (library): a Gaussian integral
of~$Q$ against a translated vacuum. The hypothesis on~$Q$ on balls and the value of the Gaussian integral with
parameter \m{(2L+2)^{-1}} give the factor \m{(2L+2)^{|\Phi_s||\ix n'||\ix n|}A\,e^{w}}.
(ii) By duality for the Gram norm (library) it suffices to bound \m{|\sigma(\lambda\cdot\operatorname{jet}^{\kappa}(F)\xi)|} for
the vectors~$\lambda$ of dual norm at most~$1$. This number is the jet of~$F$ along the polynomial
\m{z\mapsto\lambda\cdot S_\kappa(z)\xi}, to which (i) applies: the Gram norm of \m{S_\kappa(z)\xi} is at most the product
of the squared singular values of~$z$ to the row lengths times the norm of~$\xi$ (library), and
\m{x^{l}\le\exp(\log l!+lv+xe^{-v})} turns each power into the admissible growth, with \m{v=u+1/L}.
\end{proof}
\leanline{\lean{normjetle}}

\begin{proposition}[crude denominators and sizes of kernels of definite pairs]\label{p:crude}
Let \m{\ix L} and~\m{\ix N} be finite index sets, $P$ parameters, $D$ admissible data for~$P$ and \m{d\ge0} an
integer. Then there is \m{c\in\CC}, \m{c\ne0}, such that for all characters $a$, $b$ of \m{E(D)} with weights and tame
parts, every modulus~$\mathfrak{m}$, every integer \m{q\ge0} with \m{q\,g(D)^2\in\mathfrak{m}}, every function~$\Psi$
on matrices modulo~$\mathfrak{m}$ whose values are integral at every embedding into~$\CC_\ell$ with
\m{\ell\nmid q}, every polynomial~$Q$ that is homogeneous of degree~$d$, and all finite adelic points $g$, $h$, the
number \m{v=c\cdot\tk{I_{\ix L}}{I_{\ix N}}{a}{b}{\mathfrak{m}}{\Psi}{Q}(g,h)} satisfies:
\begin{enumerate}[label=\textup{(\roman*)},leftmargin=2.3em,itemsep=2pt]
\item for every prime~$p$ prime to~$q$ and every embedding \m{\iota_p\colon\CC\to\CC_p}: if $g$ and~$h$ are integral at
the primes above~$p$ and the coefficients of~$Q$ are $\iota_p$-integral, then \m{|\iota_p(v)|_p\le1};
\item if \m{|\ix L|\le40\,k(D)}, \m{|\ix N|\le40\,k(D)}, \m{d\le g(D)} and \m{q=1}, then for every \m{\sigma\in\AutC},
every integer \m{M\ge0} and every real~$B$ such that \m{|\sigma'(\Psi(x))|\le M} for all \m{\sigma'\in\AutC} and
all~$x$, and \m{|\sigma(Q(z))|\le B} for all~$z$ with \m{|\sigma(z_i)|\le g(D)^4} for all~$i$:
\m{|\sigma(v)|\le e^{x_1}\,M\,B}.
\end{enumerate}
The formal text lists a point~$h$ and a character~$b$ among the arguments of this statement; they are those
quantified inside it.
\end{proposition}
\begin{proof}
The constant~$c$ is the product of the scales that make \m{c} times the kernel equal to the jet along~$Q$ of the
scaled Taylor series of \m{\omega_{a,b}(g,h)} applied to a multiple of the datum of~$\Psi$; it depends on $D$, $d$
and the index sets only, because $Q$ is homogeneous of degree~$d$. (i) By the unit \emph{Exchange}, a power of
\m{g(D)} times the datum of~$\Psi$ is a standard datum at~$p$, since \m{q\,g(D)^2\in\mathfrak{m}}, $\Psi$ is integral
away from~$q$ and \m{p\nmid q}; the Weil representation at points integral above~$p$ preserves standard data, and jets
of standard data along $\iota_p$-integral polynomials are $\iota_p$-integral (unit \emph{Curve}). (ii) The datum is an
admissible combination of size \m{|\OEr/\mathfrak{m}|^{|\ix L||\ix N|}M} (unit \emph{Exchange}). Apply
Proposition~\ref{p:jet}\,(i) with \m{L=d} and \m{w=0}: by homogeneity the bound for~$Q$ on the polydisc of radius
\m{g(D)^4} gives the growth hypothesis with \m{A=B}. The remaining factors, \m{|\OEr/\mathfrak{m}|\le g(D)^{4f}}, the
exponential of \m{c_P(N_1|\ix L||\ix N|+1)^2} and the power of \m{2d+2}, are at most \m{e^{x_1}} by the inequalities
between the scales (unit \emph{Data}).
\end{proof}
\leanline{\lean{existsmodKernelbound}}

\subsection*{A jet that is no more divisible than predicted}
\begin{proposition}[the depth of an indicator jet]\label{p:depth}
Let $P$ be parameters and $D$ admissible data for~$P$, with nominal weights \m{w_D(\tau)} at the infinite
places~$\tau$ of \m{E(D)}, and let \m{\mathbf{m}=(\mathbf{m}_\tau)} be integers with
\m{w_D(\tau)\le\mathbf{m}_\tau\le w_D(\tau)+x^{42}} for all~$\tau$. Then there is \m{t\in\CC}, \m{t\ne0}, with
\m{|\sigma(t)|\le\exp\bigl(f(D)\,m\,b/k(D)\bigr)} for all \m{\sigma\in\AutC}, such that for every
modulus~$\mathfrak{m}$ with ideal generated by~$\bar\pi_g$, every prime~$p$ and every embedding
\m{\iota_p\colon\CC\to\CC_p} there are an exponent~$\alpha$, a family of integers \m{\alpha(\tau,i)\ge0} for
\m{1\le i\le12}, with \m{\sum_i\alpha(\tau,i)=\mathbf{m}_\tau} for all~$\tau$, and a \m{12\times1} matrix~$x_0$
modulo~$\mathfrak{m}$ with
\[p^{-F_D(p)}\le\Bigl|\iota_p\Bigl(\frac{\kappa_0^{12}}{t}\cdot\prod_\tau\mathbf{m}_\tau!\cdot\operatorname{coeff}_\alpha\,T_{\pi_g\operatorname{per}}\bigl(\mathbf{1}_{x_0}\bigr)\Bigr)\Bigr|_p .\]
Here \m{\mathbf{1}_{x_0}} is the datum of \m{(I_1,I_{12})} given by the indicator of~$x_0$ modulo~$\mathfrak{m}$ and the
vacuum, \m{\kappa_0} and \m{\operatorname{per}} are the vacuum scale and the jet scale of \m{E(D)}, \m{\pi_g} is the
conjugate of~$\bar\pi_g$, and \m{F_D(p)} is the \kw{forced part}: the sum over~$\tau$ of the divisibility by~$p$ that a
derivative of order \m{w_D(\tau)} at~$\tau$ is sure to have (unit \emph{Curve}).
\end{proposition}
\begin{proof}
\emph{The number~$t$.} Let \m{V=\operatorname{Hom}(I_1,I_{12})}. Put
\m{t=g(D)^{24f(D)}\prod_\tau\lambda(\mathbf{m}_\tau)\cdot y_0}, where \m{\lambda(l)} is the \kw{depth loss} of the unit
\emph{Curve}, the square of a product of \m{\lceil60b/k(D)\rceil} consecutive integers ending at~$l$, and $y_0$ is a
non-zero value of the integral vacuum of~$V$ at a finite vector~$y$. The bound for the conjugates of~$t$ follows
from Lemma~\ref{l:small}\,(a), the bound for the integral vacuum, \m{\mathbf{m}_\tau\le w_D(\tau)+x^{42}} and the
inequalities between the scales.

\emph{The series.} Fix $\mathfrak{m}$, $p$ and~$\iota_p$. A tight frame for~$p$ (unit \emph{Curve}) gives vectors~$w_i$
with multiplicities \m{c_i\ge3\sum_\tau\mathbf{m}_\tau}, one of them of unit size at every conjugate. With the period
factor~$S$ of \m{E(D)} let \m{\mu_i} be the linear form \m{\tau\mapsto\iota_p(S\,\overline{\tau(w_i)})}, and let
\m{\Theta=\prod_i\vartheta_p(\mu_i)^{c_i}}, where \m{\vartheta_p} is the $p$-adic Taylor series of the theta function of
the CM curve. In the ring generated by the series of the translated derivative of the Weierstrass function at
the~$\mu_i$, choose for each exponent~$\alpha$ an exponent \m{e(\alpha)}, of degree at most that of~$\alpha$, such that
the product \m{\Phi_\alpha} of the generators to the powers \m{e(\alpha)_i} with~$\Theta$ has the largest coefficient
at~$\alpha$ among all products of an element of the ring with~$\Theta$ (library). Let \m{N(\alpha)} be the absolute value
of this coefficient. By construction every product \m{H\Theta}, $H$ in the ring, is \kw{carried by one copy of~$N$}:
its coefficient at~$\alpha$ is at most \m{N(\alpha)}. The inverse \m{\Theta^{-1}} is carried by five copies: by the
$p$-adic logarithm of the theta function (library; Kronecker's first limit formula), \m{\vartheta_p=\exp(-\Lambda)} with
$\Lambda$ of order two, so that, for four $p$-adic integers~$j_b$ with \m{\sum_bj_b^2=-2}, which exist because every
natural number is a sum of four squares, \m{\Theta^{-1}} is a unit of the ring times~$\Theta$ times the four
rescalings of~$\Theta$ by the~$j_b$. Hence every \m{H\Theta\cdot\Theta^{-1}=H} is carried by six copies.

\emph{A large split.} The claim is that \m{\mathbf{m}} is a sum of twelve exponents~$\alpha_b$ with
\[\prod_b N(\alpha_b)\ \ge\ p^{-F_D(p)-\sum_\tau v_p(\lambda(\mathbf{m}_\tau))+\sum_\tau v_p(\mathbf{m}_\tau!)} .\]
If $p$ splits in~$K$, this is the limit to Katz's divisibility (unit \emph{Curve}): for each order there is a
series of the ring with a coefficient of controlled valuation; apply it to \m{\mathbf{m}} and to~$0$, each on six
copies, at the frame vector of unit size. If $p$ does not split in~$K$, then $p$ is prime to \m{g(D)} and the period
factor is a $p$-adic unit. A Frobenius packing of the orders (unit \emph{Curve}) offers two alternatives, and a large
split from pairs gives exponents carried by six copies; in the first alternative an element of the ring with
lowest term a monomial with coefficient a power of \m{\iota_p(\bar SS\,g(D))} is built by induction on the monomial
from the frame identity, in the second the Katz bound is used again; the valuations of the factorials involved are
compared by the inequality of the packing.

\emph{From the series to a datum.} The degrees satisfy \m{3e(\alpha_b)_i\le c_i}, so the tight frame produces
(unit \emph{Curve}) a finite datum of~$V$ whose scaled Taylor series is, after~$\iota_p$, the product over~$b$ of
the series \m{\Phi_{\alpha_b}} in twelve separate sets of variables; its coefficient at the exponent~$\alpha$ assembled
from the~$\alpha_b$ is \m{\prod_bN(\alpha_b)} in absolute value. This datum is a combination with $\iota_p$-integral
coefficients of translates of the vacuum, and each translate is a finite sum of indicator data modulo~$\mathfrak{m}$
with coefficients of absolute value~$1$ (library). By the ultrametric inequality one indicator datum
\m{\mathbf{1}_{x_0}} has a coefficient at~$\alpha$ at least as large. Inserting the large split and the definition
of~$t$ gives the inequality.
\end{proof}
\leanline{\lean{depthindicatorjet}}

\subsection*{Cleared values, good projections, a monomial vector}
\begin{lemma}\label{l:cleared}
In Setting~\ref{set:data}:
\begin{enumerate}[label=\textup{(\alph*)},leftmargin=2.3em,itemsep=2pt]
\item Let $p$ be a pair of~$D$ in the sense of the unit \emph{Exchange}, a character of weight~$-c$ with an
integer~$z$, let \m{\ix n} be a finite index set for which $p$ is flat, that is \m{c_\tau=z} for all~$\tau$ and
\m{|\ix n|=2z\ge4}, and \m{\sigma\in\AutC}. Then the normalised value of~$p$ at the trivial twist, the squared norm of
the trivial lift cleared of its denominators, has \m{|\sigma(\,\cdot\,)|\le\exp(x^{34})}.
\item Let $\ell$ be a prime, \m{\iota\colon\CC\to\CC_\ell} an embedding, $\vartheta$ a twist, \m{a\in\CC}, $z$ a form of
principal level on~$Y_0(D)$ such that $a$ times its image under the chain of lifts for~$\vartheta$ is integral
at~$\iota$, and \m{s\in\mathsf{G}(D)}. Then $a$ times the image of \m{|\mathsf{G}(D)|\,e_0(D)\,s\,z} under the chain is
integral at~$\iota$.
\item Let \m{\sigma\in\AutC}, $z$ a form of principal level on~$Y_0(D)$, \m{s\in\mathsf{G}(D)} and $B$ a real number
such that every value of the image of~$z$ under the chain of lifts for the trivial twist has
\m{\|\sigma(\,\cdot\,)\|^2\le e^{B}}. Then the image~$\phi$ of \m{|\mathsf{G}(D)|\,e_0(D)\,s\,z} under the chain has
\m{\nsq{\phi}{\sigma}\le\exp(x^{99}+B)}.
\item Let \m{\ix n} be a finite index set, $X$ an $L$-value datum of \m{(E(D),\ix n)} in the sense of the unit
\emph{Exchange}, with auxiliary pair~$O_X$ and jet degrees \m{d_X(\tau)}, \m{r'\ge0} an integer,
\m{e\colon\{1,\dots,r'\}\hookrightarrow\ix n} an embedding and $\alpha$ an exponent with
\m{\sum_i\alpha(\tau,i)=d_X(\tau)} for all~$\tau$. Then there is an integer coordinate vector~$v$ with
\m{P^{O_X}(\bar v)=\bigl(\prod_\tau d_X(\tau)!/\prod\alpha!\bigr)\cdot x^{\alpha}}, the monomial being read in the rows~$e$,
such that for every field~$R$ of characteristic zero and every family~$g$ of invertible \m{\ix n\times\ix n}
matrices over~$R$ whose columns \m{e(a)} are those of the identity, the weight representation of~$g$ fixes the
vector \m{w=v}, read in~$R$.
\end{enumerate}
\end{lemma}
\begin{proof}
(a) The squared norm of the trivial lift is a class sum; each value is a matrix of jets of a vacuum datum applied
to a vector, bounded by Proposition~\ref{p:jet}\,(ii) with \m{L=1}; the number of classes and the orders of the
residue rings are bounded by the units \emph{Columns} and \m{Exchange}, and the scales absorb everything.
(b), (c) By the unit \emph{Seesaw} the chain intertwines the finite group with a right translation by integral
points, so the value of the image at~$g$ is \m{\sum_{t\in\mathsf{G}(D)}|\mathsf{G}(D)|\,e_0(t)} times values of the image
of~$z$ at translates of~$g$. The numbers \m{|\mathsf{G}(D)|\,e_0(t)} are algebraic integers (unit \emph{Trace}) and
translation by integral points keeps the gauge, which gives~(b) by the ultrametric inequality. For~(c) the sum of
their conjugates' absolute values, the order of the group and the number of classes are bounded by the scales.
(d) Let \m{v_k} be the coefficient of the monomial in the integral basis polynomial~$b_k$. The monomial is
semi-invariant of the mixed weight, hence the polynomial of a coordinate vector (Input~1 of the sister note), and
by Fock duality that vector has the coordinates \m{\alpha!\,v_k/h_\kappa}; the hook number is \m{\prod_\tau d_X(\tau)!}
(unit \emph{Exchange}). The coordinates of the weight representation of~$g$ applied to~$w$ are Fock pairings of the
basis with the monomial after the row substitution by~$g\transp$, and the monomial, which involves the rows~$e$ only,
is fixed by this substitution.
\end{proof}
\leanline{\lean{normconjintAlgPart}\sep \lean{goodProjTranslateintegral}\sep \lean{normSqAtgoodProjTranslate}\sep
\lean{existsintvector}\sep the formula for the value of the image at a point is \lean{iotagoodProjTranslate}}

\subsection*{The comparison forms}
\begin{lemma}[the adjoint of the lift on a comparison form]\label{l:adj}
In Settings \ref{set:bg} and~\ref{set:data}, let $j$ be a lift index, $\vartheta$ a twist of~$D$ and \m{\kappa\in\CC}
such that, for every choice \m{i\in\Mem{j}} and every point~$g$,
\[\fnat{j}{i}(g)=\kappa\cdot\cks{O_D(\vartheta,j)}{\mathbf{b}''_D(\vartheta,j)}{\stackOverGenScalar{\splittingV{j}{\vartheta}}{i_1}}{\phi}{\lmult{\vartheta}{j}{i}{\cdot}}(g).\]
Then for every \m{F\in\Mgood{j}} and every table~$\zeta$, with \m{u_\vartheta} as in Theorem~A,
\begin{align*}
\bar\kappa\cdot\ip{F}{u_\vartheta^{*}\,\comparisonForm{j}{\vartheta}{\zeta}}
={}&|\congruenceGroup{\splittingV{j}{1}}|\cdot g(D)^{\loss{\splittingV{j}{1}}}\cdot\overline{\lineScale{j}}\\
&\cdot\sum_{i\in\Mem{j}}\overline{\weilScalar{\glueLattice{j}{1}{i}}}\int_{[\UU(W_j(D))]}\overline{F(g)\,\fnat{j}{i}(g)}\cdot\ttheta{j}{\zeta_{(i,\cdot)}}(g)\,dg .
\end{align*}
\end{lemma}
\begin{proof}
By adjointness (unit \emph{Exchange}) the left inner product is the Petersson product of~$F$ with the descent of
the comparison form. By Proposition~\ref{p:thetaprod}\,(b), for the characters \m{b=\mathbf{b}''_D(\vartheta,j)} and
\m{b'=\upsilon^{m/2-j}}, whose product is the second character of the pair, the descent of the term \m{a=(i,\mathfrak{c})}
is the conjugate of the Weil scalar times \m{\lineScale{j}} times the sum over \m{A\in\congruenceGroup{\cdot}} of the conjugate
of the coset kernel sum of the line multipliers times the kernel \m{\Theta^{-}_{\zeta_a}(\,\cdot\,,h_{\mathfrak{c}})}, both
at~\m{g\tilde A}. By hypothesis the coset kernel sum is \m{\fnat{j}{i}/\kappa}. The form~$F$ is invariant under
\m{\tilde A}, which reduces to~$1$ at the chosen places (Remark~\ref{r:liftC}), and the integral over
\m{[\UU(W_j(D))]} is invariant under right translation, so the terms of the sum over~$A$ are equal; the integrals
converge by Input~L\,(b). The sum over~$\mathfrak{c}$ of the conjugates of the kernels is \m{\ttheta{j}{\zeta_{(i,\cdot)}}}.
The splitting, its loss, its congruence group and the glued lattices do not depend on the twist.
\end{proof}
\leanline{\lean{inneradjoint}}

\begin{lemma}[integrality and size]\label{l:intsize}
In Settings \ref{set:bg} and~\ref{set:data}, let $j$ be a lift index and $\zeta$ a table.
\begin{enumerate}[label=\textup{(\alph*)},leftmargin=2.3em,itemsep=2pt]
\item Let $\vartheta$ be a twist of~$D$, $p$ a prime prime to the norms of the primes in \m{Q(\vartheta)} and
\m{\varphi\colon\CC\to\CC_p} an embedding. If
\m{\gau{\varphi\circ\mathrm{conj}}{\kappa^{-}}{0}{h_{\mathfrak{c}}}{\overline{\zeta_a}}\le1} for every index \m{a=(i,\mathfrak{c})},
then \m{\comparisonForm{j}{\vartheta}{\zeta}} is integral at~$\varphi$.
\item For every \m{\sigma\in\AutC}:
\m{\smallHookNumber{j}\cdot\nsq{\comparisonForm{j}{1}{\zeta}}{\sigma}\le\exp\bigl(2\,r\,B_0(D)\bigr)\cdot\sum_a\|\sigma(\zeta_a)\|^2}.
\end{enumerate}
\end{lemma}
\begin{proof}
(a) By Remark~\ref{r:liftC} the comparison form is a finite sum of point forms, that is of delta forms, and the
gauge of a delta form is at most the gauge of its value at its point (library). By Lemma~\ref{l:cartan}\,(b), with
\m{M=g(D)}, which lies in the ideal of the chosen primes because \m{g(D)=\bar\pi_g\overline{\bar\pi_g}}, the gauge of the
value \m{g(D)^{\operatorname{loss}}\,(y\cdot\zeta_a)} is at most the product of the gauges of~$y$ and of
\m{\overline{\zeta_a}}, the powers of \m{|g(D)|_p} cancelling. The first is at most~$1$ by Input~L\,(c), and the second
by hypothesis. (b) The squared norm is a class sum of squared norms of sums of values of delta forms. The norm of
a Cartan product is at most the product of the norms (library), \m{\|\sigma(y)\|^2\le\exp(rB_0(D))} by Input~L\,(c)
at the trivial twist, and the orders of class stabilisers divide the chosen group order (unit \emph{Data}). The
numbers of classes at the levels involved (unit \emph{Columns}), the order of the residue ring, the hook number and
\m{g(D)^{2\operatorname{loss}}} are together at most \m{\exp(rB_0(D))} by the inequalities between the scales.
\end{proof}
\leanline{\lean{comparisonFormintegral}\sep \lean{normSqAtcomparisonForm}}

\begin{lemma}[the adjoint of the comparison map on a tested form]\label{l:tt}
In Settings \ref{set:bg} and~\ref{set:data}, let $j$ be a lift index, $\vartheta$ a twist of~$D$, and let
\m{\Xi\in\FormVx{O_D(\vartheta,j)}} lie in the image of \m{e_0(D)}, in the kernel of \m{H_{O,w}(z_O(w))} for the pair
\m{O=O_D(\vartheta,j)} and a finite place~$w$ with \m{w\ne\bar w} and $w$, $\bar w$ not places of~$O$, and in the
orthogonal complement of the image of~$u_\vartheta$. Then for every choice \m{i\in\Mem{j}} and every point~$g$, with
\m{\zeta^{\Xi}} the image of~$\Xi$ under the adjoint of \m{\comparisonForm{j}{\vartheta}{\cdot}},
\m{\ttheta{j}{\zeta^{\Xi}_{(i,\cdot)}}(g)=0}.
\end{lemma}
\begin{proof}
The form~$\Xi$ satisfies Setting~\ref{set:nonlift}, with \m{\ix N=\Bix{j}}: the pair fits its parameter (unit
\emph{Exchange}); the tame part of the pair at a prime of \m{Q(\vartheta)} is that of~$\vartheta$, which is not
trivial; the reflection sums kill the image of~\m{e_0(D)} (unit \emph{Trace}); \m{\dsc{O}(\Xi)=0} because, for every
form~$y$ of principal level, \m{\ip{y}{\dsc{O}(\Xi)}=\ip{\alift_D(\vartheta,j)y}{\Xi}} and, $\Xi$ being in the image of the
self-adjoint idempotent \m{e_0(D)}, this is \m{\ip{\alift_D(\vartheta,j)(e_0(D)y)}{\Xi}=0} by Input~V and the
orthogonality to the lifts of good forms; and the rank condition holds by Input~L\,(d). Now
\m{\ttheta{j}{\zeta^{\Xi}_{(i,\cdot)}}(g)} is the inner product of \m{\zeta^{\Xi}} with the table that has the conjugate
kernel vector at \m{(g,h_{\mathfrak{c}})} in the place \m{(i,\mathfrak{c})} and zero at the other choices, hence the inner
product of~$\Xi$ with the comparison form of this table. That comparison form is a sum, over the classes of the
level~$\Gamma_i$, of sums over~$\mathfrak{c}$ of product forms as in Corollary~\ref{c:orth}, the sum over the classes
\m{\mathfrak{c}} being a class sum at the principal level \m{\bar\pi_g^{2}}, which is free. Each inner sum pairs to zero
with~$\Xi$.
\end{proof}
\leanline{\lean{tableThetaadjoint}}

\begin{proof}[Proof of Theorem A]
Take $C$ and the numbers \m{c(\vartheta)} of Input~L\,(a), with \m{u=(1\mid0)} and
\m{w=\stackOverGenScalar{\splittingV{j}{\vartheta}}{i_1}}, so that \m{uw=i_1}: the hypothesis of Lemma~\ref{l:adj} holds with
\m{\kappa=c(\vartheta)}. Write \m{R(F,\zeta)} for the right side of Lemma~\ref{l:adj}; it does not depend
on~$\vartheta$. \emph{The numbers \m{c(\vartheta)} are not zero}: otherwise all \m{\fnat{j}{i}} vanish, and by
Input~L\,(b) every good form on~$W_j(D)$ is zero, while \m{\dim\Mgood{j}\ne0}. Put
\[\gamma(\vartheta)=\overline{c(1)/c(\vartheta)},\qquad\operatorname{gl}_\vartheta=\comparisonForm{j}{\vartheta}{\cdot},\qquad
A=u_1^{*}\circ\operatorname{gl}_1 .\]
(i) \m{\gamma(\vartheta)\ne0}, \m{\gamma(1)=1} and \m{|\gamma(\vartheta)|^2=|c(1)|^2/|c(\vartheta)|^2=\HV{j}{\vartheta}/\HV{j}{1}}.
(ii) For all~$F$ and~$\zeta$, Lemma~\ref{l:adj} for $\vartheta$ and for~$1$ gives
\m{\overline{c(\vartheta)}\ip{F}{u_\vartheta^{*}\operatorname{gl}_\vartheta\zeta}=R(F,\zeta)=\overline{c(1)}\ip{F}{A\zeta}}, so
\m{u_\vartheta^{*}\operatorname{gl}_\vartheta\zeta=\gamma(\vartheta)A\zeta}.
(iii) Let $F$ be orthogonal to the image of~$A$. For a choice~$i$ and a function~$\zeta'$ on \m{\Scl}, the table
equal to~$\zeta'$ at the indices \m{(i,\mathfrak{c})} and to zero elsewhere gives
\m{0=R(F,\zeta)}, which is a non-zero constant times
\m{\int\overline{F\,\fnat{j}{i}}\cdot\ttheta{j}{\zeta'}}: the group order, \m{g(D)}, \m{\lineScale{j}} and the Weil scalar
are not zero. By Input~L\,(b), \m{F=0}. So $A$ is surjective.
(iv) Let \m{y\in\Mgood{j}} and let $\Xi$ be as in Lemma~\ref{l:tt}. Then
\m{\ip{\operatorname{gl}_\vartheta(A^{*}y)}{\Xi}=\ip{y}{A\,\zeta^{\Xi}}}, and \m{\overline{c(1)}\ip{y}{A\zeta^{\Xi}}=R(y,\zeta^{\Xi})}
vanishes because every \m{\ttheta{j}{\zeta^{\Xi}_{(i,\cdot)}}} does. So \m{\operatorname{gl}_\vartheta(A^{*}y)} is tested.
(v) and~(vi) are Lemma~\ref{l:intsize}.
\end{proof}
\leanline{\lean{existscomparisonforms}}

\sect{Where the results are used}{s:used}
Theorem~A is the input of the height bound for a ratio of lift volumes in the note \emph{The height of a ratio of
lift volumes}: (ii) and~(iii) express the Gram determinant of the adjoint lift through the comparison forms, (iv)
makes the congruence bound of a fit family applicable, (v) and~(vi) bound it at the primes and at the conjugates,
and (i) turns the twist into a ratio of Hecke values. Proposition~\ref{p:crude} is used there for the integrality
and size of the line multipliers, Proposition~\ref{p:depth}, Lemma~\ref{l:cleared}\,(a), (d) and Proposition~\ref{p:jet}
for the heights of the Hecke value ratios, and Lemma~\ref{l:cleared}\,(b), (c) in the unit \emph{Endgame}, where
crude rows complete a family of rows to a basis.

\sect{Notes on the formal text}{s:notes}
\begin{enumerate}[label=\textup{\arabic*.},leftmargin=1.6em,itemsep=5pt]
\item \emph{Arguments that only a compiler can settle.} The arguments were read from the text. In the following
declarations this reading lists a letter that the statement binds itself or that occurs only in another
declaration of the same section, and Lean presumably does not take it: a point~$h$ in the coset kernel sum and a
point~$h$ and a character~$b$ in Proposition~\ref{p:crude}; a character~$b$ in Lemmas \ref{l:thetapoint}
and~\ref{l:thetaL}; a lattice~$L$ in Lemma~\ref{l:latexists} and Proposition~\ref{p:glue}; a twist~$\vartheta$ in
Theorem~A, where it is quantified inside. Conversely, for Proposition~\ref{p:crude} the reading does not list the
two index sets $\ix L$, $\ix N$, which its uses supply as its first two arguments. The statements above follow the
uses; none of this was confirmed by asking Lean. For this reason the coset kernel sum carries no notation made by
the build.
\item \emph{Statements cut at a local definition.} The formal statements of Propositions \ref{p:jet}
and~\ref{p:crude} and of Proposition~\ref{p:thetaprod}\,(b) contain local definitions (of $u$, $M$, $F$, of~$v$, of the
translated point); the tools of this write-up read such a statement only up to the first of them. The statements
above were written from the Lean text.
\item \emph{What is chosen.} The point \m{k_{Y,e}} of a point form, the glued lattice, the divided embedding of the
splitting, the representatives of classes and the number \m{\lineScale{j}} are chosen by the axiom of choice; the
statements hold for every choice. The glued lattice for a twist is the one for the trivial twist, transported:
this is used in Lemma~\ref{l:adj}.
\item \emph{Theorem~A is stated for twists.} The maps \m{\operatorname{gl}_\vartheta} and the number \m{\gamma(\vartheta)}
are defined for every twist; the height bound later uses them at the powers of $m$-fold twists.
\item \emph{Hypotheses not used as stated.} In Lemma~\ref{l:small}\,(b) the formal statement asks that $\ix L$ be
non-empty. In Lemma~\ref{l:cartan}\,(a) the divided splitting enters only through \m{\bar\varpi_{\mathcal{V}}}.
\item \emph{Truncated subtraction.} Differences of natural numbers in the formal text, as in the depth loss and
the exponents of~$\upsilon$, are truncated at zero or are integers as the text specifies; the note writes
\m{\upsilon^{m/2-j}} for the integer exponent \m{\lfloor m/2\rfloor-j}.
\end{enumerate}

\subhead{How this note was written.} From the dossier of the unit, whose reader calls the step claims of the long
analytic proofs a first reading and notes that no second reader had looked at them. After the dossier was frozen, the second reader of the write-up compared the readings of all statements of the unit
with the Lean text and sampled its structured proofs (her message to the writer of 26 November 2026); where
her corrections concern this note they are followed: in Lemma~\ref{l:hyp} the data are the vacuum \emph{outside} the
places of the pair (the dossier had the set reversed; the Lean definition of vanishing descents confirms her
reading), the growth hypothesis of Proposition~\ref{p:jet}\,(i) is the implication stated there, and the coefficients of
the good projector in Lemma~\ref{l:cleared} are algebraic integers after multiplication by the group order. I read the six Lean files of
these declarations completely. The proof of Proposition~\ref{p:depth} (the depth of an indicator jet) I read line by
line in the Lean text. I could follow its architecture: the number~$t$, the series~$\Theta$ and the carrying by
copies of~$N$, the four squares, the case split on the splitting of~$p$, the passage from the series to an indicator
datum. I could not reconstruct the comparisons of exponents that the formal proof closes by simplification and
linear arithmetic, nor the exact statements of the inputs from the unit \emph{Curve} beyond their docstrings; the
proof above is an account of the argument, not a verification of its estimates. The same holds for
Propositions \ref{p:jet}, \ref{p:crude} and~\ref{p:period} and for Lemma~\ref{l:intsize}\,(b). The proofs of
Lemmas \ref{l:reduce} to~\ref{l:rowfix}, of Proposition~\ref{p:glue}, of Lemma~\ref{l:adj} and of Theorem~A are complete
relative to the inputs; I rewrote them from the Lean text, and they are shorter than the formal ones. I read the
Lean statements of the three results of Input~L and the definition of ``tested''; the other inputs are stated from
the dossier's reading of their docstrings. The program check of the unit's dossier covered \UnitSteps{} leaf steps
of \UnitProofs{} proofs.

\head{Appendix A. Concordance with the Lean artifact}
{\fontsize{8.6}{11.6}\selectfont\renewcommand{\leanfont}{\ttfamily\fontsize{7.7}{10}\selectfont\color{leancol}}
\setlength{\LTleft}{0pt}\setlength{\LTright}{\fill}\setlength{\tabcolsep}{0pt}
\begin{longtable}{@{}p{30mm}p{56mm}p{43mm}r@{}}
\textbf{In this note} & \textbf{Lean declaration} & \textbf{File} & \textbf{Lines}\\\midrule\endhead
\ConcordanceRows
\end{longtable}}

\subhead{The definitions of a glued lattice and of the comparison form as they are written in the artifact.}
\quoted{IsGlued}
\quoted{comparisonForm}
The formal text of Theorem~A contains a character, the sign of the orthogonal complement, that the fonts of this
note lack; it is therefore not quoted, and its row in the table above leads to it.
In these texts {\leanfont Tables j} is the space of tables, {\leanfont D.goodForms} the good forms,
{\leanfont D.goodLift} the map~$u_\vartheta$, {\leanfont Data.heckeValueProd} the number \m{\HV{j}{\vartheta}},
{\leanfont AuxPair.IsTested} the notion ``tested'', {\leanfont weylGauge} the gauge, {\leanfont LevelForm.integralForms}
the integral forms, {\leanfont normSqAt} the number \m{\nsq{\cdot}{\sigma}} and {\leanfont weylNormSq} the squared Gram
norm.

\head{Notation}
\begingroup\fontsize{8.6}{10.9}\selectfont\raggedright
\setlength{\columnsep}{9mm}\setlength{\parskip}{0pt}
\begin{multicols}{2}
\nt{$E$, $\OEr$, $\Ohat$, $\AAA$, $\AAA_f$}{a CM field; its integers; integral finite adeles; adeles; finite adeles}
\nt{$\tau$, $w$, $\bar w$, $k_w$}{an infinite place; a finite place, its conjugate, its residue field}
\nt{$\upsilon$}{the character of the Weil data}
\nt{$\ix r,\ix n,\ix N,\ix m,\ix a,\ix l,\ix L,\ix c$, $\ix n'$, $\ix n_1$, $\ix n_2$}{finite index sets}
\nt{$\mathfrak{m}$, $S$ (level places), $\mathfrak{M}$, $\mathfrak{Q}$}{a modulus; level places; products of chosen and
twist primes}
\nt{$\UU(S)(\AAA)$, $G_\infty(s)$, $\KK{\ix n}$, $\Gamma$, $K$, $K(\varpi)$}{unitary groups; stabiliser of the standard
lattice; levels}
\nt{$\clsum{h}{\Gamma}$}{the class sum}
\nt{$\tk{S}{T}{\alpha}{\beta}{\mathfrak{m}}{\Psi}{P}(g,h)$, $\omega$, $\theta$}{the theta kernel; the Weil representation; the
theta functional}
\nt{$\kappa$, $b_i$, $b_i^{-}$, $G_\kappa$, $h_\kappa$, $S_\kappa(z)$, $\lambda_\kappa(\tau,i)$}{a weight; basis polynomials;
Gram matrix; hook number; Schur matrix; row lengths}
\nt{$\operatorname{pol}_\kappa(\xi)$, $\operatorname{pol}^{\pm}_{O,\kappa}(\xi)$, $\|\xi\|_\kappa$, $\ip{\xi}{\xi'}_\kappa$}{the polynomial of a coordinate
vector, also in letters; Gram norm and form}
\nt{$O$, $O'$, $O_X$, $W_O$, $s_O$, $e_O$, $\kappa^{\pm}_O$, $\kappa_O$, $\epsilon_O$, $\mu_O$, $\psi_O$, $\mathbf{a}_O$,
$\mathbf{b}_O$, $\Gamma_O$, $\chi_O$}{an auxiliary pair and its components}
\nt{$\FormVx{O}$, $P^{O}(\xi)$, $\Upsilon_O$, $\dsc{O}(\Xi)(g)$, $H_{O,w}(X)$, $z_O(w)$}{forms, harmonic polynomial,
test datum, descent, Hecke test}
\nt{$\varpi_O$, $\bar\varpi_{\mathcal{V}}$, $\operatorname{red}_0$, $I_e$, $\tilde A$}{generator, its class, reduction,
inclusion matrix, chosen lift}
\nt{$\mathcal{V}$, $\varepsilon$, $\phi_{\mathcal{V}}$, $\phi$, $z$}{a divided splitting, its bijection, its fixed
embedding; the embedding of~$\ix r$}
\nt{$\operatorname{Stab}(\varphi,w,e)$, $\operatorname{cd}_{\chi,\varphi,w,e}(\sigma)$}{stabiliser and column datum of the
sister note}
\nt{$\gau{\iota_p}{\kappa}{\epsilon}{g}{\xi}$, $\nsq{\phi}{\sigma}$, $\ip{\phi}{\phi'}$}{gauge; squared norm at~$\sigma$;
inner product of forms}
\nt{$p$, $\ell$, $\iota_p$, $\iota$, $\varphi$, $\sigma$, $\sigma'$, $\CC_p$}{primes; embeddings; automorphisms of~$\CC$}
\nt{$\rowFixingLevel{S}{w}{e}$, $\kernelVector{O'}{b}{\mathfrak{m}}{\Psi}{\kappa}{g}{h}{i}$}{Definition~\ref{d:point}}
\nt{$\cks{O'}{b}{w}{e}{y}(g)$, $\pointForm{h}{Y}{e}{v}$, $k_{Y,e}$}{coset kernel sum; point form; its point}
\nt{$L$, $\DividedLattice{O}$, $\gamma_L$, $\reducedInverse{L}$, $\weilScalar{L}$, $\delta$}{a divided lattice and its
objects}
\nt{$\stackOverGenScalar{\mathcal{V}}{s}$, $\congruenceGroup{\mathcal{V}}$, $\IsGlued{\mathcal{V}}{s}{L}$}{Definition~\ref{d:glue}}
\nt{$\productForm{\mathcal{V}}{L}{h}{y}(\eta)$, $y\cdot\eta$, $\loss{\mathcal{V}}$}{product form; Cartan product; loss}
\nt{$P$, $D$, $E(D)$, $m$, $k(D)$, $r$, $t(D)$, $x$, $x_1$, $b$, $g(D)$, $f(D)$}{parameters and data: field, number of
lifts, ranks, step, scales, second level, number of infinite places}
\nt{$\bar\pi_g$, $\pi_g$, $\mathsf{G}(D)$, $e_0(D)$, $B_0(D)$, $w_D(\tau)$, $F_D(p)$}{the integer of norm \m{g(D)} and its
conjugate; the finite group; the good idempotent; the budget for congruences; nominal weight; forced part}
\nt{$j$, $\vartheta$, $Q(\vartheta)$, $O_D(\vartheta,j)$, $W_j(D)$, $Y_0(D)$}{lift index; twist; its primes; the auxiliary
pair; spaces}
\nt{$\Mgood{j}$, $\alift_D(\vartheta,j)$, $u_\vartheta$}{good forms; the adjoint lift; its restriction}
\nt{$\Bix{j}$, $\Six$, $\Cix$, $\Scl$, $\mathfrak{c}$, $h_{\mathfrak{c}}$, $\Mem{j}$, $i$, $i_1$}{index sets; classes and
representatives; choices}
\nt{$\mathbf{b}''_D(\vartheta,j)$, $\mathbf{a}$, $\lmult{\vartheta}{j}{i}{h}$, $\fnat{j}{i}$, $\ttheta{j}{\zeta}$,
$\HV{j}{\vartheta}$}{objects of the sister note on heights}
\nt{$\smallHookNumber{j}$, $\SmallFormValues{j}$, $\Tables{j}$, $\zeta$, $\zeta_a$, $a$}{hook number; table values;
tables; a table; an index}
\nt{$\splittingV{j}{\vartheta}$, $\glueLattice{j}{\vartheta}{i}$, $\lineScale{j}$, $\comparisonForm{j}{\vartheta}{\zeta}$,
$\Gamma_i$}{Definition~\ref{d:comp}}
\nt{$\gamma$, $A$, $\operatorname{gl}_\vartheta$, $c(\vartheta)$, $C$, $R(F,\zeta)$, $\zeta^{\Xi}$}{objects of Theorem~A and its
proof}
\nt{$K$ (field), $E_s$, $\Phi_s$, $N_1(s)$, $\Omega_s$, $\operatorname{per}_s$, $\kappa_0$, $c_P$, $N_0$, $d$}{objects of the
parameters}
\nt{$T_c(f)$, $\operatorname{jet}_Q(F)$, $\operatorname{jet}^{\kappa}(F)$, $M(L,w)$, $u$, $w_\kappa$, $B$}{scaled Taylor series;
jets; the bound of Proposition~\ref{p:jet}}
\nt{$\mathbf{m}$, $t$, $\alpha$, $x_0$, $\mathbf{1}_{x_0}$, $V$, $y_0$, $\lambda(l)$, $v_p$}{objects of
Proposition~\ref{p:depth}; $v_p$ the $p$-adic valuation}
\nt{$\Theta$, $\vartheta_p$, $\Lambda$, $\mu_i$, $w_i$, $c_i$, $S$ (period factor), $N(\alpha)$, $e(\alpha)$,
$\Phi_\alpha$, $H$, $j_b$}{objects of its proof}
\nt{$\Xi$, $\phi$, $\phi'$, $F$, $y$, $\eta$, $v$, $\xi$, $\xi'$, $\zeta'$, $\varphi_i$, $\delta_i$}{forms, functions,
coordinate vectors}
\nt{$g$, $h$, $h'$, $h_1$, $h_2$, $x$, $k$, $\gamma'$, $T$, $X$, $Y$, $Z$, $U$, $D$ (matrix), $G$, $C$, $\Delta$,
$\Delta'$}{points and matrices}
\nt{$a$, $b$, $b'$, $\alpha$, $\beta$, $e$, $e'$, $e_a$, $e_b$, $\psi$, $\psi'$, $\chi$}{characters, weights, tame parts;
$e$, $e'$ also embeddings}
\nt{$\Psi$, $\Psi_{Y,e}$, $\Psi_{L,e}$, $\mathbf{1}_{I}$, $\Theta^{\pm}$, $w_s$, $y_i(h)$}{functions on matrices; kernels;
the scaled multiplier}
\nt{$l_0$, $l$, $q$, $M$, $R$, $A$ (bound), $s$, $z$ (point), $c$, $d$, $\nu$, $\beta$ (bound), $f$, $\rho$,
$Q$}{numbers, parameters, auxiliary objects of the statements}
\nt{$\mathbb{H}$, $\mathcal{Z}$, $d_X(\tau)$, $r'$, $R$ (field), $\mathbb{Q}$}{the hyperbolic plane; a
subspace; jet degrees; an integer; a field; the rationals}
\nt{$L$ (integer, also in ``$L$-function'')}{an integer parameter of Proposition~\ref{p:jet}}
\end{multicols}
\endgroup

\end{document}
